\providecommand{\reg}[1]{\ensuremath{\mathsf{#1}}}
\section{Local assumptions and proof dependency map}\label{app:formal}
The main theorem's seed has scaled value embeddings with minimum singular value $1$, and every scaled weight/embedding block has operator norm at most $\max(4\sqrt K,\sqrt{2|\Sigma|})$. The MLP biases enforce a single local ReLU pattern. This section states the broader local assumptions used only for the outcome argument.
\begin{assumption}[Small initialization]\label{ass:init}
Write $\theta_0=\varepsilon\tilde\theta$. Every scaled weight and embedding block has operator norm at most $a_0$, and the value embedding block has smallest singular value at least $c_{\mathrm{emb}}>0$. Each scaled MLP-bias coordinate has absolute value at least $1$, and each bias vector has norm at most $N_b$. Require $\varepsilon\le c_{\mathrm{init}}\min(1,N_b^{-1})$, for a sufficiently small constant depending on $a_0,c_{\mathrm{emb}},K$ and the numbers of blocks and heads. Put $r_0=\min(1/8,c_{\mathrm{emb}}/4)$ and
\[
\mathcal B_\varepsilon=\{\theta:\|\theta-\theta_0\|\le r_0\varepsilon,\ \beta=0\}.
\]
All parameter norms are Euclidean/Frobenius unless the tied metric of \cref{eq:metric} is explicitly invoked.
\end{assumption}
\begin{assumption}[Task]\label{ass:task}
Use the independent-letter modular-counting task in \cref{def:task}. The outcome proof requires $0<\lambda<1$.
\end{assumption}
The enlarged ball has radius $2r_0\varepsilon$; bias coordinates there have magnitude at least $3\varepsilon/4$, and value-embedding singular values are at least $\varepsilon c_{\mathrm{emb}}/2$. For sufficiently small $c_{\mathrm{init}}$, MLP input contributions are $O(\varepsilon^2)$ and cannot change the ReLU signs. A trained relative bias is treated separately in \cref{lem:bias-uniform}; it changes the input-independent base attention distribution, not the Fourier support count.

\paragraph{Credit notation.}\label{def:credit}
For $\xi=e_y-u$, $u=\mathbf1/K$, put $C_{\mathrm{out}}(\theta)=\E\langle z_\theta,\xi\rangle$. Its gradient is the answer-dependent part of the cross-entropy gradient. A term in a score expansion reads the input letters appearing in its monomial. This support count is local to each term, not to the complete network.

\paragraph{Dependencies.}
\Cref{prop:represent} uses only the explicit finite-score construction in Appendix~\ref{app:architectures}. For Theorem~\ref{thm:matched}(i), \cref{lem:target-contraction} and the architecture estimates \cref{prop:network,prop:score-degree} imply \cref{lem:score-sum}; \cref{lem:bias-uniform,lem:end-smooth} then justify the local loss floor and the gradient-flow escape argument. Appendix~\ref{app:match} transfers its metric to the common tied family. For part (ii), \cref{lem:J-features,lem:J-symmetry,lem:J-reduce} derive the invariant parameter dynamics; \cref{lem:J-geometry,lem:J-levels,lem:J-gradient,lem:J-drift,lem:J-align,lem:J-floor,lem:J-grow,lem:J-descent} prove the aligned bootstrap, time bound, margins and persistence. Part (iii) intersects the intervals and substitutes the width/scale regime. Numerical checks test identities and finite executors; they do not replace these proofs.

\section{Exact two-block executors}
\label{app:architectures}

This section proves \cref{prop:represent}. An implementation of the construction is in the accompanying code repository.

\paragraph{Architecture.}
\label{app:construction}
\label{app:model}
Each block has three heads with $d_{\mathrm h}=2n+6$ (the number of positions) and one ReLU MLP of width $\max\{|\Sigma|,3K(n+1),2n+1\}$;
there is no LayerNorm:
\begin{align*}
 \bm h_i&\leftarrow \bm h_i+\sum_{h=1}^{3}\sum_{j\le i}\operatorname{softmax}_{j\le i}\Bigl(\tfrac{1}{\sqrt{d_{\mathrm h}}}(\mathbf W^{(h)}_Q\bm h_i)^{\!\top}\mathbf W^{(h)}_K\bm h_j\Bigr)\mathbf W^{(h)}_V\bm h_j,\\
 \bm h_i&\leftarrow \bm h_i+\mathbf W_{\mathrm{out}}\relu(\mathbf W_{\mathrm{in}}\bm h_i+\bm b_{\mathrm{in}}),\qquad \mathrm{logits}_i=\mathbf V\,\bm h_i .
\end{align*}

\paragraph{Registers.} The residual stream is the concatenation of one-hot registers
\[
 \underbrace{\reg{tl}}_{|\Sigma|}\ \underbrace{\reg{tv}}_{K}\ \underbrace{\reg{pos}}_{d_{\mathrm h}}\ \underbrace{\reg{qry}}_{|\Sigma|}\ \underbrace{\reg{match}}_{1}\
 \underbrace{\reg{prev}}_{K}\ \underbrace{\reg{mt}}_{1}\ \underbrace{\reg{cnt}}_{1}\ \underbrace{\reg{s0}}_{K}\ \underbrace{\reg{out}}_{K},
 \qquad d_{\mathrm{res}}=2|\Sigma|+4K+d_{\mathrm h}+3 .
\]
$\reg r[a]$ is coordinate $a$ of register $\reg r$, and $\mathbf E_{\reg a\leftarrow\reg b}$ copies $\reg b$ into $\reg a$.

\paragraph{Layout.} Token ids: letters $0,\ldots,|\Sigma|-1$, value $v\mapsto|\Sigma|+v$, then SEP, ``:'', EOS. Positions ($0$-indexed):
\[
\begin{array}{ll}
 \text{prompt} & s_0@0,\quad w^\star@1,\quad w_t@(1+t),\quad \mathrm{SEP}@(n+2)\\
 \text{step row} & s_t@(n+2+t),\quad ``{:}"@(2n+3),\quad s_n@(2n+4)\\
 \text{outcome row} & ``{:}"@(n+3),\quad s_n@(n+4)
\end{array}
\]
Position $i$ predicts the token at $i+1$: $n+1+t$ predicts $s_t$, and $n+3$ predicts $s_n$ in the outcome row.

\paragraph{Programs.} State arithmetic is modulo $K$.
\begin{center}\small
\fitwidth{\begin{tabular}{@{}lll@{}}
\toprule
 & step executor & outcome executor\\
\midrule
block 1, head 1 & \multicolumn{2}{l}{$\reg{qry}\leftarrow$ the query, at every letter}\\
block 1, MLP & \multicolumn{2}{l}{$\reg{match}\leftarrow b_t$}\\
block 2, head 1 & $\reg{prev}\leftarrow\widetilde s_{t-1}$ & $\reg{cnt}\leftarrow$ mean of $\reg{match}$ over the letters\\
block 2, head 2 & $\reg{mt}\leftarrow\reg{match}$ of letter $t$ & $\reg{s0}\leftarrow s_0$\\
block 2, head 3 & copy the last displayed count to the answer position & ---\\
block 2, MLP & $\reg{out}\leftarrow\mathrm{lookup}(\reg{prev},\reg{mt})$ & $\reg{out}\leftarrow\mathrm{lookup}(\reg{s0},\reg{cnt})$\\
\bottomrule
\end{tabular}}
\end{center}

\paragraph{Route and average.} Let $\Omega_n=\max\{16,3\log(8(n+3))\}$, $C_Q=\Omega_n\sqrt{d_{\mathrm h}}$, and let $\mathrm{src}(i)\subseteq\{0,\ldots,i\}$ be the source set
of position $i$ (a fallback $\{f\}$ where none is listed). Set
\begin{equation}
 \begin{gathered}
 (\mathbf W_K)_{j,\reg{pos}[j]}=1,\qquad (\mathbf W_Q)_{j,\reg{pos}[i]}=C_Q\ \ (j\in\mathrm{src}(i)),\\
 \tfrac{1}{\sqrt{d_{\mathrm h}}}(\mathbf W_Q\bm h_i)^\top \mathbf W_K\bm h_j=\Omega_n\,[j\in\mathrm{src}(i)] .
 \end{gathered}
 \label{eq:route}
\end{equation}
With $\mathbf W_V=\mathbf E_{\reg a\leftarrow\reg b}$, the head writes into $\reg a$ the average of $\reg b$ over $\mathrm{src}(i)$, up to the leak bounded below; $|\mathrm{src}(i)|=1$ is a
route and $|\mathrm{src}(i)|>1$ an average. A fallback $f>i$ gives uniform attention at a position with no supervised prediction.

\paragraph{Explicit source sets and embeddings.}
Token embeddings write the letter one-hot into $\reg{tl}$ and count one-hot into $\reg{tv}$; format embeddings are zero. Position embeddings write $e_i$ into $\reg{pos}$. All other initial registers vanish. Block 1 head 1 routes $1+t$ to $1$ for $1\le t\le n$ (fallback $0$), writing $\reg{tl}$ into $\reg{qry}$. Its other heads are zero. In the step executor, block 2 head 1 routes $n+2$ to $0$ and $n+1+t$ to itself for $2\le t\le n$, with fallback $1$, writing $\reg{tv}$ into $\reg{prev}$. Head 2 routes $n+1+t$ to $1+t$ for $1\le t\le n$, with fallback $0$, writing $\reg{match}$ into $\reg{mt}$. Head 3 routes $2n+3$ to $2n+2$, with fallback $1$, writing $\reg{tv}$ into $\reg{out}$. In the outcome executor, head 1 routes $n+3$ to $\{2,\ldots,n+1\}$ (fallback $0$), writing $n\reg{match}$ into $\reg{cnt}$; head 2 routes $n+3$ to $0$ (fallback $1$), writing $\reg{tv}$ into $\reg{s0}$; head 3 is zero. The step readout puts format weight $20$ on positions $2n+2$ for colon and $2n+4$ for EOS; the outcome readout uses positions $n+2$ and $n+4$. All masks are full causal masks. Unused tensor entries are zero.

\paragraph{Lookup.} For a count $s$ and a bit $b\in\{0,1\}$, use one unit per pair $(s',b')$ with
\[
 \mathrm{pre}_{s'b'}=[s=s']+[b=b']-\tfrac{3}{2}=\begin{cases}\tfrac{1}{2}&(s',b')=(s,b),\\\le-\tfrac{1}{2}&\text{otherwise},\end{cases}
\]
and output weights $2\bm e_{s'+b'}$:
\begin{equation}
 \textstyle\sum_{s',b'}2\,\bm e_{s'+b'}\relu(\mathrm{pre}_{s'b'})=\bm e_{s+b}.
 \label{eq:lookup}
\end{equation}

\subsection{Both programs are exact}
\label{app:exact}

\begin{proof}[Proof of \cref{prop:represent}]
\emph{Ideal routing} ($\Omega_n=\infty$). With $m=\sum_tb_t$ and $[m=a]=\relu(m-a+1)-2\relu(m-a)+\relu(m-a-1)$,
\begin{gather*}
 \reg{match}_{1+t}=2\textstyle\sum_{a\in\Sigma}\relu\bigl([w_t=a]+[w^\star=a]-\frac{3}{2}\bigr)=b_t,\\
 \reg{prev}_{n+1+t}=\bm e_{s_{t-1}},\qquad \reg{mt}_{n+1+t}=b_t,\qquad
 \reg{cnt}_{n+3}=n\cdot\tfrac{1}{n}\textstyle\sum_tb_t=m,\qquad \reg{s0}_{n+3}=\bm e_{s_0},\\
 \reg{out}_{n+1+t}=\sum_{j,\iota}\bm e_{j+\iota}\,2\relu\Bigl([s_{t-1}=j]+[b_t=\iota]-\frac{3}{2}\Bigr)=\bm e_{s_t},\\
 \reg{out}_{n+3}=\sum_{j,a}\bm e_{j+a}[s_0=j][m=a]=\bm e_{s_n}.
\end{gather*}
The test $[s_0=j]$ is applied by adding $(n+2)([s_0=j]-1)$ to each of the three pre-activations of $[m=a]$. The readout has weight $20$
on $\reg{out}$ for value tokens and weight $20$ on the relevant $\reg{pos}$ coordinates for ``:'' and EOS; all unused maps are zero.

\emph{Finite scores.} Put $\omega=(d_{\mathrm h}-1)e^{-\Omega_n}$. Every intended source set at a required position is nonempty and causal, so its outside softmax mass is at most $\omega<1/[256(n+3)^2]$. The intended source weights are equal. Coordinatewise,
\[\begin{gathered}
|\reg{match}-b_t|\le2\omega,\quad
\|\reg{prev}-e_{s_{t-1}}\|_\infty\le\omega,\quad
|\reg{mt}-b_t|\le3\omega,\\
|\reg{cnt}-m|\le3n\omega,\quad
\|\reg{s0}-e_{s_0}\|_\infty\le\omega.
\end{gathered}\]
The infinity norm here is intentional: an unqualified Euclidean bound with the same constant would be false for a leaking one-hot copy. The matching step lookup unit has preactivation at least $1/2-4\omega$; every other unit has preactivation at most $-1/2+4\omega<0$. It thus contributes at least $1-8\omega$ to the correct output and zero to every wrong output. The additional head-3 write at these positions is nonnegative with each coordinate at most $\omega$ (its fallback is the query letter, which has zero count register). The resulting count margin is at least $20(1-9\omega)>19$.

For the outcome lookup, let $T_a(x)=\relu(x-a+1)-2\relu(x-a)+\relu(x-a-1)=\max(1-|x-a|,0)$. For the true start-state row its argument is
$x=\reg{cnt}+(n+2)(\reg{s0}[s_0]-1)$, so
$|x-m|\le(4n+2)\omega=:D<1/192$. Consequently $T_m(x)\ge1-D$, while the sum of all other integer tents with $0\le a\le n$ is at most $D$: for $|x-m|<1$, only the adjacent tent can be nonzero. For a wrong start-state row, $x\le n+3n\omega-(n+2)(1-\omega)<-1$, so all its tents vanish. Modular collisions can only add nonnegative output to a class. The correct output is therefore at least $1-D$ and every wrong output at most $D$, giving margin at least $20(1-2D)>19$.

At the step final-answer position head 3 copies the last displayed count, with correct coordinate at least $1-\omega$ and wrong coordinates at most $\omega$. The fallback previous-state head has $\|\reg{prev}\|_\infty\le\omega$, so its lookup is inactive there. The final-answer margin is at least $20(1-2\omega)>19$. At colon and EOS prediction positions, the position readout gives the gold format token logit $20$, the previous-state lookup is inactive, and any spurious count-copy output is at most $\omega$. Thus format margins are also positive. The same argument covers outcome colon/EOS positions, where the count lookup's start-state fallback has norm at most $\omega$ and its tents vanish. This verifies every required prediction, including the first step whose source is $s_0$ in the prompt.

\emph{Open neighbourhood.} Let $\mathcal P_n$ be the finite set of gold prefixes and supervised positions. By continuity there is
$r_{\mathrm{exec},n}>0$ with
\[
 \|\bm\theta-\bm\theta^\star\|<r_{\mathrm{exec},n}\implies\max_{P\in\mathcal P_n}\|\bm z_{\bm\theta}(P)-\bm z_{\bm\theta^\star}(P)\|_\infty<\tfrac{19}{2},
\]
so every greedy prediction is $a^\star$, and induction over the continuation gives the claim.
\end{proof}



\section{Fourier credit and the outcome plateau}
\label{app:outcome}
\subsection{Fourier contraction}
\label{app:C-fourier}
\label{app:assumptions}




(A1)--(A3) denote \cref{ass:arch,ass:init,ass:task}, and (A3$'$) a uniform start. Unless stated otherwise the query and the start
count are fixed. Constants $c,C$ depend only on $a_0,\embconst,K$ and the numbers of blocks and heads. Every bound stated on $\initball$ also
holds on $\initball'=\{\bm\theta:\|\bm\theta-\bm\theta_0\|\le2\ballradius\initscale,\ \beta=0\}$, whose relative radius is at most $\min(\frac{1}{4},\frac{\embconst}{2})$; there
$|b_m|\ge(1-2\ballradius)\initscale\ge\frac{3}{4}\initscale$ and $\sigma_{\min}(\mathbf E_{\mathrm{val}})\ge\frac{\initscale\embconst}{2}$. The only width-dependent threshold is $\initscale_0=c\min(1,\frac{1}{N_b})$.

\begin{lemma}[Gradient split]
\label{lem:gradient-split}
For every differentiable model, with $\outcredit(\bm\theta)=\E[\langle\bm z_{\bm\theta},\bm y-\bm u\rangle]$,
\[
 \nabla_{\bm\theta}\Lout=\underbrace{\E\bigl[(\Jz)^\top(\bm p-\bm u)\bigr]}_{\text{answer-independent}}-\underbrace{\nabla_{\bm\theta}\outcredit}_{\text{answer-dependent}},\qquad
 \nabla_{\bm\theta}\outcredit=\E\bigl[(\Jz)^\top(\bm y-\bm u)\bigr].
\]
\end{lemma}

\begin{proof}
The logit gradient of the cross-entropy is $\bm p-\bm e_y=(\bm p-\bm u)-(\bm y-\bm u)$.
\end{proof}

\paragraph{Conditioning on unseen letters.}
\label{app:target-spectrum}

\paragraph{Contraction.}\label{app:contraction} By \cref{eq:lambda}, $\mixrate^2=1-2p(1-p)(1-\cos\tfrac{2\pi}{K})$, and for $K=2$,
$\mixrate=|1-2p|$ and $1-\mixrate^2=4p(1-p)$. With $\alpha=e^{2\pi\mathrm i/K}$ and the Fourier multipliers
$\rho_j=1-p+p\alpha^j$, so that $\mixrate=\max_{j\ne0}|\rho_j|$, one unseen letter multiplies each nonconstant count
mode $\alpha^{js_t}$ by $\rho_j$, while one displayed transition keeps the innovation
\begin{equation}
 \E\bigl[\bigl|\alpha^{j s_t}-\rho_j\,\alpha^{j s_{t-1}}\bigr|^2\bigr]=1-|\rho_j|^2\ge1-\mixrate^2 .
 \label{eq:move-variance}
\end{equation}
In the Fourier basis of $\one^\perp$, $\operatorname{Cov}(\bm g_t)$ is diagonal with entries $\tfrac{1}{K}(1-|\rho_j|^2)=\tfrac{2p(1-p)}{K}\bigl(1-\cos\tfrac{2\pi j}{K}\bigr)$,
which bounds the covariance below by $(1-\lambda^2)/K$ on $\mathbf1^\perp$.

\begin{lemma}[Conditional contraction]
\label{lem:target-contraction}
Assume (A3) and fix the start count and the query, so that the letters are independent. Let $\bm\xi=\bm y-\bm u$, $S\subseteq[n]$ and $\mathcal F_S=\sigma(w_t:t\in S)$. Then
\begin{equation}
 \bigl\|\E[\bm\xi\mid\mathcal F_S]\bigr\|^2=\frac{1}{K}\sum_{j=1}^{K-1}|\rho_j|^{2(n-|S|)}\le\Bigl(1-\frac{1}{K}\Bigr)\mixrate^{2(n-|S|)} .
 \label{eq:horizon}
\end{equation}
Hence for every square-integrable $\mathcal F_S$-measurable matrix $\mathbf H$ with $K$ rows,
 $\bigl\|\E[\mathbf H^\top\bm\xi]\bigr\|\le\mixrate^{n-|S|}\,\|\mathbf H\|_{L^2},\qquad \|\mathbf H\|_{L^2}:=\bigl(\E\|\mathbf H\|_F^2\bigr)^{1/2}.$ 
\end{lemma}

\begin{proof}
Write $s_n=s_0+\sum_tb_t$. The letters are independent, so for $j\ne0$
\[
 \E\bigl[\alpha^{js_n}\bigm|\mathcal F_S\bigr]=\alpha^{js_0}\prod_{t\in S}\alpha^{jb_t}\prod_{t\notin S}\E\bigl[\alpha^{jb_t}\bigr]=\alpha^{js_0}\,\rho_j^{\,n-|S|}\prod_{t\in S}\alpha^{jb_t},
\]
whose modulus is $|\rho_j|^{n-|S|}$. Parseval over the nonconstant modes $j=1,\ldots,K-1$ gives the identity, and $|\rho_j|\le\mixrate$ the bound.
The matrix bound follows from \cref{eq:horizon} by the tower property
and Cauchy--Schwarz.
\end{proof}

\paragraph{Sharpness.}\label{rem:tight} The $j=1$ term of \cref{eq:horizon} is $\frac{1}{K}\mixrate^{2(n-|S|)}$, and $\mathbf H=\E[\bm\xi\mid\mathcal F_S]$ attains the second bound up to the
factor $(1-\frac{1}{K})^{1/2}$, so the exponent $n-|S|$ is sharp for $0<\mixrate<1$.

\subsection{Transformer score expansion}
\label{app:model-spectrum}

The contraction estimate helps only if we can control how many letters a network term uses. Expanding the attention scores does
this without asking the entire output to depend on a small set: each term has its own set of letters, and terms with larger sets
are smaller near initialization.

\paragraph{Score expansion.} For $\zeta\in[0,1]$, let $\bm z(\zeta)$ be the network with every score $\psi_{il}$, in both blocks, replaced by
$\zeta\psi_{il}$. Then $\bm z(1)=\bm z_{\bm\theta}$, and for every integer $k_0\ge0$,
\begin{equation}
 \bm z_{\bm\theta}=\sum_{k=0}^{k_0}\bm z^{[k]}+\bm z_{\mathrm{rem}},\qquad
 \bm z^{[k]}=\frac{1}{k!}\frac{d^k\bm z}{d\zeta^k}(0),\qquad
 \bm z_{\mathrm{rem}}=\int_0^1\frac{(1-\zeta)^{k_0}}{k_0!}\frac{d^{k_0+1}\bm z}{d\zeta^{k_0+1}}(\zeta)\,d\zeta .
 \label{eq:score-split}
\end{equation}
$\mathbf J^{[k]}$ and $\mathbf J_{\mathrm{rem}}$ are the corresponding Jacobian terms. A term \emph{reads} the letters whose inputs it contains.

\begin{proposition}[Score count and letter support]
\label{prop:score-degree}
Fix the start state and the query. For $\bm\theta\in\initball$, $\Jz=\sum_{k=0}^{k_0}\mathbf J^{[k]}+\mathbf J_{\mathrm{rem}}$, and there is $\attconst$ such that:
\begin{enumerate}[label=(\roman*),leftmargin=*,itemsep=2pt,topsep=2pt]
\item $\mathbf J^{[0]}=\mathbf J^{[0]}_\emptyset+\sum_{t=1}^n\mathbf J^{[0]}_t$, where $\mathbf J^{[0]}_\emptyset$ reads no letter, $\mathbf J^{[0]}_t$ reads only $w_t$, $\sup\|\mathbf J^{[0]}_\emptyset\|_F\le C\initscale$ and
$\sum_t\sup\|\mathbf J^{[0]}_t\|_F\le C\initscale^2$;
\item for $k\ge1$, $\mathbf J^{[k]}=\sum_\nu\mathbf T_{k,\nu}$, where each $\mathbf T_{k,\nu}$ reads at most $2k+1$ letters and $\sum_\nu\sup\|\mathbf T_{k,\nu}\|_F\le C\initscale^2(\attconst\initscale^4)^k$;
\item $\sup\|\mathbf J_{\mathrm{rem}}\|_F\le C\initscale^2(\attconst\initscale^4)^{k_0+1}$.
\end{enumerate}
The same bounds hold for $\bm z^{[k]}$ and $\bm z_{\mathrm{rem}}$ with one more factor $\initscale$. Suprema are over inputs; constants do not depend on
$n$, $k$ or $k_0$.
\end{proposition}

$\attconst$ is fixed once, large enough for \cref{prop:score-degree,prop:network}; it depends only on $a_0$, $\embconst$ and $K$.

\label{app:network}

\begin{proposition}[Bounds for the score expansion]
\label{prop:network}
Assume (A1)--(A2). Uniformly on $\initball$, at every position and input, with $\bm x_i$ the token-plus-position vector,
all derivatives with respect to $s$ are evaluated at $s=0$ along $\bm\theta+sv$.
Let $\|F\|_{\mathrm{abs}}$ be the absolute sum of $F$ at $\zeta=0$: expand $F$ into monomials, products of
parameter blocks, inputs $\bm x_j$ and input-independent attention weights; bound each monomial by the product of the norm bounds of its
parameter factors on the ball, of $\sup\|\bm x_j\|$ for its inputs, and of its weights; and sum these bounds before any cancellation:
\begin{enumerate}[label=(\roman*),leftmargin=*,itemsep=3pt,topsep=4pt]
\item with input-independent offsets $\bm O^1=\bm o_1$, $\bm O^2=\bm o_1+\bm o_2$: $\bm h_i^\ell=\bm x_i+\bm O^\ell+O(\initscale^2)$, $|\psi_{il}|\le C\initscale^4$,
$\|\bm z_i\|\le C\initscale^2$, $\|\dd{\bm z_i}{s}\|\le C\initscale\|v\|$;
\item for $k\ge1$, $\zeta\in[0,1]$ and $\|v\|=1$,
\[
 \frac{1}{k!}\Bigl\|\frac{d^k\bm z_i}{d\zeta^k}\Bigr\|\le C\initscale^3(C_0\initscale^4)^k,\qquad
 \frac{1}{k!}\Bigl\|\dd{}{s}\left(\frac{\mathrm d^k\bm z_i}{\mathrm d\zeta^k}\right)\Bigr\|\le C\initscale^2(C_0\initscale^4)^k;
\]
the same holds for the absolute sums at $\zeta=0$, for $\bm h^2_i$ with one power of $\initscale$ fewer, and for the Taylor remainder with $k=k_0+1$;
\item the zero-score network is affine in the inputs; its transported logit terms have tangent $O(\initscale^2)$, and direct-input and
input-free terms have tangent $O(\initscale)$; an order-$k$ coefficient contains at most $2k+1$ inputs; a one-score logit term has size
$O(\initscale^7)$ and tangent $O(\initscale^6)$;
\end{enumerate}
The constants do not depend on $n$, the widths or $k$.
\end{proposition}

\begin{proof}
For a network quantity $F$ let $F^{[k]}(\zeta)=\frac{1}{k!}\frac{d^kF}{d\zeta^k}(\zeta)$, so that $\bm z^{[k]}=\bm z^{[k]}(0)$, and write $\|F^{[k]}\|$ for its supremum over
inputs, positions and $\zeta\in[0,1]$ (the $\ell_1$ norm for attention rows).

\emph{(i).} $|b_m|\ge\frac{3}{4}\initscale>C\initscale^2\ge|(\mathbf W_{\mathrm{in}}\bm h)_m|$, so one activation matrix $\mathbf D_{\mathrm{act}}$ serves every input and every $\zeta$, and
\[
 \begin{aligned}
 \bm h_i^\ell&=\mathbf M_\ell\bigl(\bm h_i^{\ell-1}+\bm a_i^\ell\bigr)+\bm o_\ell,\\
 \mathbf M_\ell&=\mathbf I+\mathbf W^\ell_{\mathrm{out}}\mathbf D_{\mathrm{act}}\mathbf W^\ell_{\mathrm{in}},\qquad
 \bm o_\ell=\mathbf W^\ell_{\mathrm{out}}\mathbf D_{\mathrm{act}}\bm b^\ell.
 \end{aligned}
\]
with $\bm a_i^\ell=\sum_{\mathrm{heads}}\sum_{j\le i}\omega_{ij}\mathbf W_V\bm h^{\ell-1}_j$ and $\omega_{i\cdot}=\softmax(\zeta\psi_{i\cdot})$. Then
 \[
\begin{gathered}
\|\mathbf M_\ell-\mathbf I\|\le C\initscale^2,\\ \|\bm o_\ell\|\le CN_b\initscale^2\le C\initscale,\\ \|\dd{\bm o_\ell}{s}\|\le(a_0+2N_b)\initscale\le C,
\end{gathered}
\] 
and $\|\omega_{i\cdot}\|_1=1$. Induction over the two blocks gives (i).

\emph{Softmax calculus.} Let $\omega=\softmax(s)$ and
$\E_\omega[f]=\sum_j\omega_jf_j$. For a direction $\bm\nu$,
\begin{align*}
 \left.\dd{}{t}\omega_j(s+t\bm\nu)\right|_{t=0}
 &=\omega_j(\nu_j-\E_\omega[\bm\nu]),\\
 \left.\dd{}{t}\E_{\omega(s+t\bm\nu)}[f]\right|_{t=0}
 &=\E_\omega[\bm\nu f]-\E_\omega[\bm\nu]\E_\omega[f].
\end{align*}
After $k$ derivatives, each term has at most $k+1$ differentiable factors
and contains each direction once. The absolute coefficient sums satisfy
$a_{k+1}\le2(k+1)a_k$, $a_0=1$. Summing against the probability weights gives
\begin{equation}
 \bigl\|D^k\omega(s)[\bm\nu_1,\ldots,\bm\nu_k]\bigr\|_1\le2^kk!\prod_a\|\bm\nu_a\|_\infty .
 \label{eq:softmax-avg}
\end{equation}
For the smooth finite-dimensional maps in the fixed-ReLU region, the
higher-order chain rule \citep[Theorem~2]{clark2013faa} yields
\[
 \frac{1}{j!}\|D^jF\|\le Ac^j,\quad \|h^{[k]}\|\le ab^k
 \quad\Longrightarrow\quad
 \|(F\circ h)^{[k]}\|\le Aca\,b^k(1+ca)^{k-1},\quad k\ge1.
\]
Here the inner map's derivative bounds hold for every positive order; the
composition count is $\binom{k-1}{j-1}$.

\emph{(ii), values.} With $\|\bm\omega^{\ell,[0]}\|_1=1$, $\|\bm h^{0,[k]}\|=0$ for $k\ge1$, and $a_r\ge1$ in the compositions,
\begin{alignat*}{2}
 \|\psi^{\ell,[k]}\|&\le C\initscale^2\textstyle\sum_{a+b=k}\|\bm h^{\ell-1,[a]}\|\,\|\bm h^{\ell-1,[b]}\| &\qquad& \why{score is a product}\\
 \|(\zeta\psi^\ell)^{[k]}\|&\le\|\psi^{\ell,[k]}\|+\|\psi^{\ell,[k-1]}\| && \why{$(\zeta\psi)^{[k]}=\zeta\psi^{[k]}+\psi^{[k-1]}$}\\
 \|\bm\omega^{\ell,[k]}\|_1&\le\textstyle\sum_{j=1}^k2^j\sum_{a_1+\cdots+a_j=k}\prod_r\|(\zeta\psi^\ell)^{[a_r]}\| && \why{composition, \cref{eq:softmax-avg}}\\
 \|\bm a^{\ell,[k]}\|&\le C\initscale\textstyle\sum_{a+c=k}\|\bm\omega^{\ell,[a]}\|_1\,\|\bm h^{\ell-1,[c]}\| && \why{head output is a product}\\
 \|\bm h^{\ell,[k]}\|&\le C\bigl(\|\bm h^{\ell-1,[k]}\|+\|\bm a^{\ell,[k]}\|\bigr) && \why{affine block}\\
 \|\bm z^{[k]}\|&\le C\initscale\,\|\bm h^{2,[k]}\| && \why{readout}
\end{alignat*}
With $b=C\initscale^4$, block $1$ gives, for $k\ge1$, $\|\bm\omega^{1,[k]}\|_1\le b^k$ and $\|\bm a^{1,[k]}\|+\|\bm h^{1,[k]}\|\le C\initscale^2b^k$ (while $\|\bm h^{1,[0]}\|\le C\initscale$). Then $\|\psi^{2,[0]}\|\le C\initscale^4$,
$\|\psi^{2,[k]}\|\le C\initscale^5(k+1)b^k$ for $k\ge1$, and $\|(\zeta\psi^2)^{[k]}\|\le C(2b)^k$, so
\begin{gather*}
 \|\bm\omega^{2,[k]}\|_1\le C\sum_{j=1}^k\binom{k-1}{j-1}C^j(2b)^k\le C(Cb)^k,\\
 \|\bm a^{2,[k]}\|+\|\bm h^{2,[k]}\|\le C\initscale^2(Cb)^k,\qquad \|\bm z^{[k]}\|\le C\initscale^3(C_0\initscale^4)^k .
\end{gather*}

\emph{(ii), tangents, remainders, absolute sums.} Differentiating along $v$ replaces one factor of each product by its tangent, or one
$(\zeta\psi)^{[a]}$ in a softmax term, or adds $D^{j+1}\omega[\dd{(\zeta\psi)^{[0]}}{s},\ldots]$. With $\|\dd{\bm h^{0,[0]}}{s}\|=O(1)$, $\|\dd{\psi^{\ell,[0]}}{s}\|=O(\initscale^3)$
and $j+1\le2^k$ absorbed into $C_0^k$,
\begin{gather*}
 \|\dd{\bm\omega^{\ell,[k]}}{s}\|_1\le C\initscale^{-1}(C_0\initscale^4)^k,\qquad \|\dd{\bm a^{\ell,[k]}}{s}\|+\|\dd{\bm h^{\ell,[k]}}{s}\|\le C\initscale(C_0\initscale^4)^k,\\
 \|\dd{\bm z^{[k]}}{s}\|\le C\initscale^2(C_0\initscale^4)^k .
\end{gather*}
The integral remainder is bounded by the supremum of the next normalized derivative.

\emph{Absolute sums.} $\|\cdot\|_{\mathrm{abs}}$ is subadditive and submultiplicative, and $\|F\|\le\|F\|_{\mathrm{abs}}$. At $\zeta=0$ every attention row is an
input-independent probability vector $\bm\pi$ (uniform over the visible sites here, $\softmax(\beta)$ in \cref{lem:bias-uniform}), and
 $\Bigl\|\textstyle\sum_j\pi_j\prod_rF_r(j)\Bigr\|_{\mathrm{abs}}\le\prod_r\max_j\|F_r(j)\|_{\mathrm{abs}} .$ 
At $\zeta=0$, $D^j\omega$ expands as $\omega(0)$-weights times products of direction entries and $\omega(0)$-averages, with total coefficient at most $2^jj!$,
so \cref{eq:softmax-avg} holds for absolute sums with $\|\bm\nu_a\|_\infty$ replaced by $\max_l\|\nu_{a,l}\|_{\mathrm{abs}}$. Hence every line of the recursion
above holds verbatim for $\|\cdot\|_{\mathrm{abs}}$ at $\zeta=0$, which gives the absolute-sum bounds of (ii).

\emph{Tangents of absolute sums.} Take the norm bound of every parameter factor at least $\initscale$ (weights and embeddings $(a_0+1)\initscale$,
biases $(N_b+1)\initscale$). An order-$k$ monomial has at most $C(k+1)$ factors: at most $k$ scores, each with $\mathbf W_Q$ and $\mathbf W_K$, and at
most $2k+1$ inputs, each carried by at most four parameter blocks. Replacing one factor by its unit-direction tangent multiplies the bound by at
most $\frac{1}{\initscale}$, so the tangent bound is at most $\frac{C(k+1)}{\initscale}\le\frac{C2^k}{\initscale}$ times the monomial bound, uniformly in the unit direction
$v$; this lets $\sup_v$ pass inside the sum over monomials.

\emph{(iii).} At zero scores the network is affine in the inputs (for a head with a submask, $\frac{1}{m_j}$ becomes one over the size of its
visible set):
\begin{gather*}
 \bm h^1_j=\mathbf M_1\Bigl(\bm x_j+\frac{1}{m_j}\sum_{l\le j}\mathbf W^1_V\bm x_l\Bigr)+\bm o_1,\\
 \bm z^{[0]}_i=\mathbf V\Bigl[\mathbf M_2\Bigl(\bm h^1_i+\frac{1}{m_i}\sum_{j\le i}\mathbf W^2_V\bm h^1_j\Bigr)+\bm o_2\Bigr].
\end{gather*}
Transported terms contain $\mathbf V$, $\mathbf W_V$ and some $\bm x_j$, so their tangent is $O(\initscale^2)$; direct-input and input-free terms have tangent $O(\initscale)$.
With $d(F)$ the number of input occurrences in $F$,
\begin{alignat*}{2}
 d(\bm h^{\ell,[a]})&\le2a+1 &\qquad& \why{induction on $\ell$}\\
 d(\psi^{\ell,[b]})&\le2b+2 && \why{score is a product of two residuals}\\
 d\bigl((\zeta\psi)^{[a]}\bigr)&\le2a && \why{$(\zeta\psi)^{[a]}=\psi^{[a-1]}$}\\
 d(\text{order-$k$ head term})&\le2a_1+\cdots+2a_j+2c+1=2k+1 && \why{$a_1+\cdots+a_j+c=k$}
\end{alignat*}
and affine MLPs and parameter derivatives add none. The one-score correction is
 $\frac{1}{m_i}\sum_{j\le i}\bigl(\psi_{ij}(0)-\bar\psi_i(0)\bigr)\mathbf W_V\bm h_j(0),\qquad \bar\psi_i(0)=\frac{1}{m_i}\sum_{l\le i}\psi_{il}(0);$ 
its logit terms contain $\mathbf V,\mathbf W_V,\mathbf W_Q,\mathbf W_K$ and three residual inputs, so they have size $O(\initscale^7)$ and tangent $O(\initscale^6)$.

\end{proof}

\begin{proof}[Proof of \cref{prop:score-degree}]
Expand each coefficient into monomials in parameters and inputs. A Jacobian entry of a monomial is a monomial with the same inputs, and the
Frobenius norm of a map into $\R^K$ is at most $\sqrt K$ times its operator norm.

\noindent(i) At zero scores the network is affine in the inputs (\cref{prop:network}(iii)): each term contains at most one input. Group the terms by
that input. At the answer position the direct input is fixed and contains no letter; the transported terms have tangent $O(\initscale^2)$, and
their coefficients sum absolutely to a constant because every zero-score attention row is a probability vector.

\noindent(ii) By \cref{prop:network}(iii) an order-$k$ monomial contains at most $2k+1$ inputs, and by \cref{prop:network}(ii) the absolute sum of the
tangents at $\zeta=0$ is at most $C\initscale^2(C_0\initscale^4)^k$; take $\attconst\ge C_0$.

\noindent(iii) Apply \cref{prop:network}(ii) with $k=k_0+1$. Logit bounds carry one more factor $\initscale$ than Jacobian bounds.
\end{proof}

\subsection{Exponential outcome plateau}
\label{app:thm1}

Combining \cref{lem:target-contraction,prop:network} gives the following credit bound.

\begin{lemma}[Credit summation]
\label{lem:score-sum}
Assume (A1)--(A3), $\bm\theta\in\initball$, $0<\initscale\le\initscale_0$ and $0<\mixrate<1$, with the start state and query fixed. Let
$k_0=\lceil\frac{n-1}{2}\rceil$. If $\attconst\initscale^4\le\frac{1}{2}\mixrate^2$, then
 $\bigl\|\E\bigl[\mathbf J_{\mathrm{score}}^\top(\bm y-\bm u)\bigr]\bigr\|\le C\initscale^6\mixrate^{n-3},\qquad \mathbf J_{\mathrm{score}}=\Jz-\mathbf J^{[0]}.$ 
If moreover $\initscale\le\mixrate$, then
\begin{equation}
 \|\nabla_{\bm\theta}\outcredit(\bm\theta)\|\le C\initscale\mixrate^n,\qquad |\outcredit(\bm\theta)|\le C\initscale^2\mixrate^n.
 \label{eq:thm1-full}
\end{equation}
The bounds persist after averaging over the start state and query.
\end{lemma}

\begin{proof}
Apply \cref{lem:target-contraction} to each term of \cref{prop:score-degree}, with $\|\mathbf H\|_{L^2}\le\sup\|\mathbf H\|_F$ and $\|\bm y-\bm u\|\le1$ for the remainder:
\begin{alignat*}{2}
 \bigl\|\E\bigl[(\mathbf J^{[0]})^\top(\bm y-\bm u)\bigr]\bigr\|&\le C\initscale\mixrate^n+C\initscale^2\mixrate^{n-1} &\qquad& \why{(i): $0$ or $1$ letter}\\
 \bigl\|\E\bigl[(\mathbf J^{[k]})^\top(\bm y-\bm u)\bigr]\bigr\|&\le C\initscale^2(\attconst\initscale^4)^k\mixrate^{n-2k-1} && \why{(ii): $\le2k+1$ letters, $k\ge1$}\\
 &=C\initscale^6\mixrate^{n-3}\Bigl(\frac{\attconst\initscale^4}{\mixrate^2}\Bigr)^{k-1} && \\
 \bigl\|\E\bigl[\mathbf J_{\mathrm{rem}}^\top(\bm y-\bm u)\bigr]\bigr\|&\le C\initscale^2(\attconst\initscale^4)^{k_0+1} && \why{(iii)}\\
 &=C\attconst\initscale^6\mixrate^{2k_0}\Bigl(\frac{\attconst\initscale^4}{\mixrate^2}\Bigr)^{k_0} && \\
 &\le C\initscale^6\mixrate^{n-3}\Bigl(\frac{\attconst\initscale^4}{\mixrate^2}\Bigr)^{k_0} && \why{$2k_0\ge n-1$}
\end{alignat*}
If $2k+1>n$, the factor $\mixrate^{n-2k-1}\ge1$ still bounds the credit. Since $\frac{\attconst\initscale^4}{\mixrate^2}\le\frac{1}{2}$, the
score-dependent terms sum to at most $3C\initscale^6\mixrate^{n-3}$, which is the first bound. Adding the zero-order term,
\begin{align*}
 \|\nabla_{\bm\theta}\outcredit\|&\le C\bigl(\initscale\mixrate^n+\initscale^2\mixrate^{n-1}+\initscale^6\mixrate^{n-3}\bigr),\\
 |\outcredit|&\le C\bigl(\initscale^2\mixrate^n+\initscale^3\mixrate^{n-1}+\initscale^7\mixrate^{n-3}\bigr),
\end{align*}
by \cref{lem:gradient-split} and the same argument for the logits, and $\initscale\le\mixrate$ reduces both to \cref{eq:thm1-full}.
\end{proof}

For $c$ small enough, $\initscale\le c\mixrate$ implies $\initscale\le\mixrate$ and $\attconst\initscale^4\le\frac{1}{2}\mixrate^2$, and $\initscale^2\le c\mixrate$ implies $\attconst\initscale^4\le\frac{1}{2}\mixrate^2$.



\paragraph{A trained relative bias.} Let $\bm\vartheta$ denote the parameters other than $\beta$, and
$\mathcal B^{\mathrm{full}}_{\initscale}=\{\bm\theta:\|\bm\vartheta-\bm\vartheta_0\|\le2\ballradius\initscale\}$ with $\beta$ unrestricted. A zero readout in $\tilde{\bm\theta}$ is admissible in (A2).

\begin{lemma}[Uniformity in the relative bias]
\label{lem:bias-uniform}
Assume (A1)--(A3), $\initscale\le\initscale_0$. On $\mathcal B^{\mathrm{full}}_{\initscale}$, with constants independent of $\beta$:
\begin{enumerate}[label=(\alph*),leftmargin=*,itemsep=3pt,topsep=4pt]
\item the logit bounds of \cref{prop:network} and of \cref{prop:score-degree} for $\bm z^{[k]},\bm z_{\mathrm{rem}}$ hold;
\item if $\initscale\le c\mixrate$, then $|\outcredit(\bm\theta)|\le C_2\initscale^2\mixrate^n$;
\item the conclusions of \cref{lem:end-smooth} hold for unit directions in all parameters, $\beta$ included.
\end{enumerate}
\end{lemma}

\begin{proof}
Each attention row is
 $\bm\omega_{i\cdot}(\zeta)=\softmax\bigl(\beta_{i-\cdot}+\zeta\psi_{i\cdot}\bigr)\qquad\text{over its visible sites.}$ 

\noindent(a) At $\zeta=0$,
 $\bm\omega_{i\cdot}(0)=\softmax(\beta_{i-\cdot}),\qquad \bm\omega_{i\cdot}(0)\ge0,\qquad \|\bm\omega_{i\cdot}(0)\|_1=1,$ 
which does not depend on the input; the zero-score network is affine with the same input counts $d(\cdot)$. Since $\beta$ does not
depend on $\zeta$,
 $(\beta+\zeta\psi)^{[k]}=\zeta\psi^{[k]}+\psi^{[k-1]}\qquad(k\ge1),$ 
so the recursion of \cref{prop:network}(ii) is unchanged, and \cref{eq:softmax-avg} holds uniformly in $\beta$. The ReLU pattern stays
fixed, since $|(\mathbf W_{\mathrm{in}}\bm h)_m|\le C\initscale^2<|b_m|$ for any attention weights. The uniform weights $\frac{1}{m_j}$ enter only
\cref{prop:network}, which is not needed here.

\smallskip\noindent(b) The proof of \cref{lem:score-sum} uses only \cref{lem:target-contraction} and the logit bounds of (a), so
$|\outcredit(\bm\theta)|\le C_2\initscale^2\mixrate^n$.

\smallskip\noindent(c) A unit direction with $\beta$-component $v_\beta$ perturbs each row's scores by $\bm\nu=v_\beta$, $\|\bm\nu\|_\infty\le1$. By
\cref{eq:softmax-avg} with $k=1,2$, the first two row derivatives have $\ell_1$ norms at most $C$. The product-rule logit bounds in \cref{lem:end-smooth} therefore apply unchanged.
\end{proof}

\begin{lemma}[Smoothness near initialization]
\label{lem:end-smooth}
Assume (A1)--(A2), $\initscale\le\initscale_0$. There is $C_1$, independent of $n$, $\initscale$ and the widths, such that for every $\bm\theta\in\initball'$ and
every unit direction $\bm v$,
 $\|\nabla_{\bm\theta}\Lout(\bm\theta)\|\le C_1\initscale,\qquad \Bigl|\frac{d^2}{ds^2}\Lout(\bm\theta+s\bm v)\big|_{s=0}\Bigr|\le C_1 .$ 
\end{lemma}

\begin{proof}
Along $\bm\theta+s\bm v$ at $s=0$, applying the product rule twice to the
two blocks and \cref{eq:softmax-avg} with $k=1,2$ gives
\[
 \|\bm z\|\le C\initscale^2,\qquad
 \left\|\dd{\bm z}{s}\right\|\le C\initscale,\qquad
 \left\|\frac{\mathrm d^2\bm z}{\mathrm d s^2}\right\|\le C.
\]
The ReLU pattern is fixed by \cref{prop:network}(i). These bounds use only
operator norms and attention-row $\ell_1$ norms bounded by one, so $C$ is
independent of $n$ and the widths.
For cross-entropy, the log-sum-exp derivative formula
\citep[Section~3.1.5]{boyd2004convex} gives
\[
 \nabla_{\bm z}\ell=\bm p-\bm e_y,\qquad
 \nabla_{\bm z}^2\ell=\operatorname{diag}\bm p-\bm p\bm p^\top,
 \qquad \|\nabla_{\bm z}\ell\|\le\sqrt2,\quad
 \|\nabla_{\bm z}^2\ell\|\le1.
\]
The chain rule therefore yields
$|\dd{\ell}{s}|\le C\initscale$ and
$|\mathrm d^2\ell/\mathrm d s^2|\le C$.
Average over prompts and take the supremum over unit $\bm v$.
\end{proof}

\paragraph{Outcome plateau.}\label{app:end} For one prompt, $\KL(\bm u\|\bm p)=\log\sum_ae^{z_a}-\log K-\langle\bm u,\bm z\rangle$, hence
\begin{equation}
 \Lout(\bm\theta)=\log K+\E\bigl[\KL(\bm u\,\|\,\bm p_{\bm\theta})\bigr]-\outcredit(\bm\theta).
 \label{eq:lo-identity}
\end{equation}

\paragraph{Ambient projected flow.}
For this supporting estimate, let $P(t)$ be an orthogonal projection onto trainable ambient parameter blocks, with finitely many switches on each bounded interval, and let $\dot\theta=-P(t)\nabla_\theta\Lout$ start at the zero-readout seed. The calculation below gives the local loss, displacement and assigned-probability bounds until time $r_0^2\lambda^{-n}/C_2$. Full ambient joint flow is included. Appendix~\ref{app:match} supplies the separate metric comparison required for the tied scalar flow in Theorem~\ref{thm:matched}.

\begin{proof}[Ambient outcome plateau]
\emph{Loss floor.} Let $B_n=C_2\initscale^2\mixrate^n$, $\initscale\le c\mixrate$ and $\bm\theta\in\mathcal B^{\mathrm{full}}_{\initscale}$.
\begin{alignat}{2}
 \Lout(\bm\theta)&=\log K+\E\,\KL(\bm u\,\|\,\bm p_{\bm\theta})-\outcredit(\bm\theta) &\qquad& \why{\cref{eq:lo-identity}}\nonumber\\
 &\ge\log K-B_n && \why{$\KL\ge0$, \cref{lem:bias-uniform}(b)}\label{eq:lo-floor}
\end{alignat}

\emph{Escape time.} By \cref{lem:end-smooth} and \cref{lem:bias-uniform}(c), $\Lout$ is smooth on $\mathcal B^{\mathrm{full}}_{\initscale}$, so the flow is defined between the
switching times of $\mathbf P(t)$ and continuous across them. Let $T_{\mathrm{exit}}=\inf\{t:\|\bm\theta(t)-\bm\theta_0\|>\ballradius\initscale\}$, the norm
including $\beta$. For $t\le T_{\mathrm{exit}}$,
\begin{alignat*}{2}
 \ddt{\Lout}&=-\|\mathbf P\nabla_{\bm\theta}\Lout\|^2=-\Bigl\|\ddt{\bm\theta}\Bigr\|^2 &\qquad& \why{$\mathbf P$ constant between switches}\\
 \int_0^t\Bigl\|\ddt{\bm\theta}\Bigr\|^2&=\Lout(\bm\theta_0)-\Lout(\bm\theta(t))\le B_n && \why{$\Lout(\bm\theta_0)=\log K$, \cref{eq:lo-floor}}\\
 \|\bm\theta(t)-\bm\theta_0\|^2&\le t\int_0^t\Bigl\|\ddt{\bm\theta}\Bigr\|^2\le t\,B_n && \why{Cauchy--Schwarz}\\
 T_{\mathrm{exit}}&\ge\frac{\ballradius^2}{C_2}\,\mixrate^{-n} && \why{$\|\bm\theta(T_{\mathrm{exit}})-\bm\theta_0\|=\ballradius\initscale$}
\end{alignat*}
and \cref{eq:lo-floor} holds on $[0,T_{\mathrm{exit}}]$.

\emph{Near-chance prediction.}
\begin{alignat*}{2}
 \E\,\KL(\bm u\,\|\,\bm p_t)&\le\outcredit(\bm\theta(t))\le B_n &\qquad& \why{\cref{eq:lo-identity}, $\Lout(\bm\theta(t))\le\log K$}\\
 \E[\bm p_{t,y}]-\tfrac{1}{K}&\le\E\,\mathrm{TV}(\bm p_t,\bm u)\le\sqrt{\frac{C_2}{2}}\;\initscale\mixrate^{n/2} && \why{Pinsker, Jensen}
\end{alignat*}

\emph{Constants.} Take $c$ the smaller of $\frac{\ballradius^2}{C_2}$ and the constant for which $\initscale\le c\mixrate$ gives $\initscale\le\mixrate$ and
$\attconst\initscale^4\le\frac{1}{2}\mixrate^2$, and $C=\max(C_2,\sqrt{\frac{C_2}{2}})$.
\end{proof}

\section{Process learning in the structured Transformer}
\label{app:joint}
\label{app:step-supplied-routing}



\subsection{Model and exact features}

We count modulo $K$ with a fixed query of probability $p\in(0,1)$ and a uniform start; every non-query letter has the same embedding. Let $\bm x_t=\bm e_{s_{t-1}}$, $\mathbf R\bm e_s=\bm e_{s+1}$, $r_b=\sigma\bigl(-(-1)^ba\bigr)$ for $b\in\{0,1\}$ and $\sigma(z)=(1+e^{-z})^{-1}$.

\paragraph{Seed.} Fix orthonormal residual directions $\hat{\bm e}_0,\hat{\bm e}_\ell$, a count register $(\hat{\bm e}^{\mathrm c}_s)_{s\in\mathbb Z_K}$, a gate register
$(\hat{\bm e}^{\mathrm g}_s)_{s\in\mathbb Z_K}$, and a unit vector $\hat{\bm e}_{\mathrm h}$ of the head space. The embeddings are
\[
 E(s)=\initscale\,\hat{\bm e}^{\mathrm c}_s,\qquad E(\text{query letter})=\initscale(\hat{\bm e}_0-\hat{\bm e}_\ell),\qquad E(\text{other letters})=\initscale(\hat{\bm e}_0+\hat{\bm e}_\ell).
\]
The gate head is head $1$ of block $2$. At the position $i_t$ that holds $s_{t-1}$ and predicts $s_t$ ($t\ge2$), its fixed mask admits
$i_t$ and the position $j_t$ of the current letter, and its relative bias is zero. Its maps are $\mathbf W_Q=q\,\bm\Pi_Q$, $\mathbf W_K=k\,\bm\Pi_K$ and
$\mathbf W_V=v\,\bm\Pi_V$, with scalars $q,k,v$ and fixed matrices
\[
 \bm\Pi_Q=\textstyle\sum_s\hat{\bm e}_{\mathrm h}(\hat{\bm e}^{\mathrm c}_s)^{\!\top},\qquad \bm\Pi_K=\hat{\bm e}_{\mathrm h}\hat{\bm e}_\ell^{\top},\qquad \bm\Pi_V=\sum_s\hat{\bm e}^{\mathrm g}_s(\hat{\bm e}^{\mathrm c}_s)^{\!\top}.
\]
All other embeddings, heads and MLP weights are zero, and the MLP biases are $\pm\initscale$. The only trainable readout blocks are the
$K\times K$ blocks $\mathbf U$ (count register) and $\mathbf W=w(\mathbf I-\mathbf R)$ (gate register), both initialized at zero.
The query and key train along $\bm\Pi_Q,\bm\Pi_K$; the value map $v\bm\Pi_V$ is fixed (\cref{rem:J-line}).

\paragraph{Training.} The trained parameters are $\bm\vartheta=(\mathbf U,\mathbf W,q,k)$, with $\mathbf W$ on the line $\R(\mathbf I-\mathbf R)$, and
$\ddt{\bm\vartheta}=-\nabla_{\bm\vartheta}\Lstep(\bm\vartheta)$ for the Euclidean (Frobenius) norm on $\R^{K\times K}\times\R(\mathbf I-\mathbf R)\times\R\times\R$. The value factor $v$
and all other parameters are fixed. The four trained parameters follow this one flow at all times; there are no stages and no stopping
rule. The initialization is
 $\mathbf U(0)=0,\qquad w(0)=0,\qquad \tfrac{1}{2}\initscale\le|q(0)|,\ |k(0)|,\ |v|\le2\initscale,$ 
with arbitrary signs.

\begin{lemma}[Exact features]
\label{lem:J-features}
For every $\bm\vartheta$ and $t\ge2$, the logits of the prediction of $s_t$ are
 $\bm z_t=\initscale \mathbf U\bm x_t+\initscale v\,r_{b_t}\,\mathbf W\bm x_t,\qquad a=\frac{\initscale^2qk}{\sqrt{d_{\mathrm h}}}.$ 
\end{lemma}

\begin{proof}
Block $1$ and all MLPs add zero, so each residual is its embedding. The query at $i_t$, the keys and the values are
\begin{gather*}
 \text{query}_{i_t}=\initscale q\,\hat{\bm e}_{\mathrm h},\qquad \text{key}_{j_t}=(-1)^{b_t}\initscale k\,\hat{\bm e}_{\mathrm h},\qquad \text{key}_{i_t}=0,\\
 \text{value}_{i_t}=\initscale v\,\hat{\bm e}^{\mathrm g}_{s_{t-1}},\qquad \text{value}_{j_t}=0 .
\end{gather*}
The scores are $(-1)^{b_t}a$ at $j_t$ and $0$ at $i_t$, so $i_t$ gets weight $r_{b_t}$. The residual at $i_t$ becomes
 $\initscale\hat{\bm e}^{\mathrm c}_{s_{t-1}}+\initscale v\,r_{b_t}\,\hat{\bm e}^{\mathrm g}_{s_{t-1}},$ 
and the readout gives $\bm z_t$.
\end{proof}

\begin{lemma}[Cyclic symmetry]
\label{lem:J-symmetry}
$\mathbf R\mathbf U(t)=\mathbf U(t)\mathbf R$ for all $t\ge0$; $\mathbf W(t)\in\R(\mathbf I-\mathbf R)$ commutes with $\mathbf R$ by definition.
\end{lemma}

\begin{proof}
Shifting $s_0$ by one preserves the data law and the match bits, and $\mathrm{CE}(\bm z,j+1)=\mathrm{CE}(\mathbf R^{-1}\bm z,j)$ for the cross-entropy. By \cref{lem:J-features},
 $\Lstep\circ S=\Lstep,\qquad S(\mathbf U,\mathbf W,q,k)=(\mathbf R^{-1}\mathbf U\mathbf R,\ \mathbf W,\ q,\ k).$ 
$S$ is a linear isometry, so $S\bm\vartheta(t)$ solves the flow, with $S\bm\vartheta(0)=\bm\vartheta(0)$. By uniqueness, $S\bm\vartheta(t)=\bm\vartheta(t)$.
\end{proof}

\subsection{Reduction}

Let $N=K-2$, $\rho=\frac{2}{K}$, $\pi_0=1-p$, $\pi_1=p$, $H_p=-p\log p-(1-p)\log(1-p)$ and $p_{\min}=\min(p,1-p)$. For $m,s\in\R$ let
 $F(m;s)=\log\bigl(1+e^{-m}+N e^{-s-m/2}\bigr),\qquad \phi(m;s)=-\pd{F}{m}(m;s);$ 
for $K=2$, $F(m;s)=\ell(m):=\log(1+e^{-m})$.

\begin{lemma}[Reduced loss and flow]
\label{lem:J-reduce}
Write $\mathbf U=\sum_ju_j\mathbf R^j$. Then $u_j=u_{\mathrm{bg}}$ for $2\le j\le K-1$. Let
 $c=\initscale(u_0-u_1),\qquad s=\initscale\bigl(\tfrac{1}{2}(u_0+u_1)-u_{\mathrm{bg}}\bigr),\qquad g=2\initscale vw,\qquad m_b=(-1)^b\bigl(c+g\,r_b\bigr),$ 
with $s=0$ if $K=2$. The prediction of $s_t$ with match bit $b_t=b$ puts logits $s+\frac{m_b}{2}$ on its target, $s-\frac{m_b}{2}$ on the other count of
$\{s_{t-1},s_{t-1}+1\}$, and $0$ on the remaining $N$ counts, up to a common shift. Hence, for every $n$,
 $\Lstep=\mathcal L(c,g,s,a):=(1-p)\,F(m_0;s)+p\,F(m_1;s).$ 
With $\qkbalance=q(0)^2-k(0)^2$,
\begin{equation}
 \begin{gathered}
 \ddt{c}=-\rho\initscale^2\pd{\mathcal L}{c},\qquad \ddt{g}=-\rho\initscale^2v^2\pd{\mathcal L}{g},\qquad \ddt{s}=-\frac{\initscale^2}{2N}\pd{\mathcal L}{s}\ (K\ge3),\qquad \ddt{a}=-\frac{\initscale^4S}{d_{\mathrm h}}\pd{\mathcal L}{a},\\
 S:=q^2+k^2=\sqrt{\qkbalance^2+\frac{4d_{\mathrm h}a^2}{\initscale^4}}.
 \end{gathered}
 \label{eq:J-flow}
\end{equation}
\end{lemma}

\begin{proof}
\emph{Background coordinates.} By \cref{lem:J-features}, $\Lstep$ depends on $\mathbf U$ through $(u_j)$, and the targets are at offsets $0$ and $1$.
Permuting $u_2,\ldots,u_{K-1}$ preserves $\Lstep$ and $\|\mathbf U\|_F^2=K\sum_ju_j^2$ and fixes $\mathbf U(0)=0$; as in \cref{lem:J-symmetry},
$u_j=u_{\mathrm{bg}}$ for $j\ge2$.

\emph{Logits.} For $\bm x_t=\bm e_{s'}$, $\mathbf W\bm x_t=w(\bm e_{s'}-\bm e_{s'+1})$. By \cref{lem:J-features}, after subtracting $\initscale u_{\mathrm{bg}}$,
 $z_{t,s'}=s+\tfrac{1}{2}(c+g\,r_{b_t}),\qquad z_{t,s'+1}=s-\tfrac{1}{2}(c+g\,r_{b_t}),\qquad z_{t,s'+j}=0\ \ (2\le j\le K-1).$ 
The target is $s'+b_t$, so its logit is $s+\frac{m_{b_t}}{2}$ and the other one is $s-\frac{m_{b_t}}{2}$. The cross-entropy is
$\log(1+e^{-m}+Ne^{-s-m/2})=F(m;s)$.

\emph{Flow.} At a circulant $\mathbf U$ the Frobenius gradient $\nabla_{\mathbf U}\Lstep$ is circulant (it is invariant under the isometry $S$ of
\cref{lem:J-symmetry}), so the flow stays on circulant matrices and acts on their coefficients as $\ddt{u_j}=-\frac{1}{K}\pd{\Lstep}{u_j}$ and $\ddt{w}=-\frac{1}{2K}\pd{\Lstep}{w}$, since
$\|\mathbf U\|_F^2=K\sum_ju_j^2$ and $\|\mathbf W\|_F^2=2Kw^2$. Substituting the definitions of $c,g,s$ in \cref{lem:J-reduce}
into these metric gradients gives their equations in \cref{eq:J-flow}.
For the gate,
\begin{gather*}
 \ddt{q}=-\frac{\initscale^2k}{\sqrt{d_{\mathrm h}}}\pd{\mathcal L}{a},\qquad \ddt{k}=-\frac{\initscale^2q}{\sqrt{d_{\mathrm h}}}\pd{\mathcal L}{a},\qquad \frac{d}{dt}(q^2-k^2)=0.
\end{gather*}
Differentiating $a=\initscale^2qk/\sqrt{d_{\mathrm h}}$ gives the remaining equation of \cref{eq:J-flow}; $S^2=(q^2-k^2)^2+4q^2k^2$ gives its expression for $S$. The flow exists for all $t\ge0$: the vector field is smooth, $|\pd{\mathcal L}{c}|,|\pd{\mathcal L}{g}|,|\pd{\mathcal L}{s}|\le1$, so
$c,g,s$ grow at most linearly, and $|\pd{\mathcal L}{a}|\le\frac{|g|}{4}$, so $(q,k)$ grows at most exponentially by Gr\"onwall.
\end{proof}

\paragraph{Notation.} $r_0+r_1=1$ and $\tau:=r_1-r_0=\tanh(\frac{a}{2})$. Let
\begin{gather*}
 \mathrm{err}_b=\pi_b\,\phi(m_b;s),\qquad \mathrm{err}=\mathrm{err}_0+\mathrm{err}_1,\\
 h=\tfrac{1}{2}(m_0-m_1)=c+\tfrac{g}{2},\qquad R=m_0+m_1=-g\tau,\qquad Q=g-\tfrac{v^2}{2}c,
\end{gather*}
the weighted errors, the count-only margin, the rule contrast from \cref{eq:rule}, and a gate--count balance. Let $P_b$ be the
probability that the prediction with match bit $b$ assigns to its target.

\begin{lemma}[Derivatives]
\label{lem:J-geometry}
\begin{enumerate}[label=(\alph*),leftmargin=*,itemsep=2pt,topsep=3pt]
\item $\pd{\mathcal L}{c}=\mathrm{err}_1-\mathrm{err}_0$, $\quad\pd{\mathcal L}{g}=r_1\mathrm{err}_1-r_0\mathrm{err}_0=\tfrac{1}{2}\bigl(\pd{\mathcal L}{c}+\tau\,\mathrm{err}\bigr)$, $\quad\pd{\mathcal L}{a}=g\,\dd{\sigma(a)}{a}\,\mathrm{err}$.
\item $0<\phi<1$; $\phi(m;s)\ge\frac{1}{2}\sigma(-m)$; $\phi(m;s)\ge\frac{1}{2}(1-P_b)$ at $m=m_b$; $m\mapsto\phi(m;s)$ is nonincreasing; and
$\phi(m-2x;s)\le e^{2x}\phi(m;s)$ for $x\ge0$.
\item $F(m;s)\ge\ell(m)\ge-m$.
\item $\ddt{\mathcal L}=-\rho\initscale^2(\pd{\mathcal L}{c})^2-\rho\initscale^2v^2(\pd{\mathcal L}{g})^2-\frac{\initscale^2}{2N}(\pd{\mathcal L}{s})^2-\frac{\initscale^4S}{d_{\mathrm h}}(\pd{\mathcal L}{a})^2\le0$ (no $s$ term if $K=2$).
\end{enumerate}
\end{lemma}

\begin{proof}
(a) Differentiate the loss and margins in \cref{lem:J-reduce}, using
$\pd{F}{m}=-\phi$ and $r_0+r_1=1$.

(b) With $A=e^{-m}$ and $B=Ne^{-s-m/2}$, $\phi=\frac{A+B/2}{1+A+B}$: the probability of the other count plus half the background
probability. Hence $0<\phi<1$ and
 $\phi\ge\frac{1}{2}\cdot\frac{A+B}{1+A+B}=\frac{1}{2}(1-P_b)\ge\frac{1}{2}\cdot\frac{A}{1+A}=\frac{1}{2}\sigma(-m).$ 
Convexity of log-sum-exp \citep[Section~3.1.5]{boyd2004convex}, composed with these affine logits, makes $F(\cdot;s)$ convex and $\phi$ nonincreasing. Finally
 $\frac{d}{dm}\log\phi=-\frac{A+B/4}{A+B/2}+\frac{A+B/2}{1+A+B}>-1,$ 
and integration over $[m-2x,m]$ gives the last claim.

(c) Drop $Ne^{-s-m/2}$; $\log(1+e^{-m})\ge-m$. (d) is the chain rule with \cref{eq:J-flow}.
\end{proof}

\begin{lemma}[Loss levels]
\label{lem:J-levels}
Let $\mathcal L^\star=\frac{1}{2}p_{\min}\log2$.
\begin{enumerate}[label=(\alph*),leftmargin=*,itemsep=2pt,topsep=3pt]
\item If $\mathcal L<p_{\min}\log2$, then $P_0,P_1>\frac{1}{2}$: every prediction of $s_t$, $t\ge2$, has positive margin against all $K-1$ other counts.
\item If $P_0,P_1\ge\frac{1}{2}$, then $\mathrm{err}_0^2+\mathrm{err}_1^2\ge \frac{\mathcal L^2}{32}$.
\item If $\mathcal L>\mathcal L^\star$, then $\mathrm{err}_0^2+\mathrm{err}_1^2\ge \frac{p_{\min}^2}{267}$.
\end{enumerate}
\end{lemma}

\begin{proof}
(a) $\pi_b\ge p_{\min}$, so $-\log P_b=F(m_b;s)<\log2$.

(b) For $P\ge\frac{1}{2}$, $-\log P\le\frac{1-P}{P}\le2(1-P)$. By \cref{lem:J-geometry}(b), $F(m_b;s)\le4\phi(m_b;s)$, so $\mathcal L\le4\,\mathrm{err}$ and
$\mathrm{err}_0^2+\mathrm{err}_1^2\ge\frac{1}{2}\mathrm{err}^2\ge\frac{1}{32}\mathcal L^2$.

(c) If $P_b<\frac{1}{2}$, then $\mathrm{err}_b\ge\frac{\pi_b}{4}\ge\frac{p_{\min}}{4}$ by \cref{lem:J-geometry}(b); otherwise (b) applies. So
 $\mathrm{err}_0^2+\mathrm{err}_1^2\ge\min\Bigl\{\frac{p_{\min}^2}{16},\ \frac{(\mathcal L^\star)^2}{32}\Bigr\}=\frac{(\log2)^2}{128}\,p_{\min}^2\ge\frac{p_{\min}^2}{267}. \qedhere$ 
\end{proof}

\begin{lemma}[Gradient lower bound]
\label{lem:J-gradient}
$(\pd{\mathcal L}{c})^2+(\pd{\mathcal L}{g})^2\ge\frac{1}{3}\tau^2(\mathrm{err}_0^2+\mathrm{err}_1^2)$. If $v^2\le1$, then
 $\ddt{\mathcal L}\le-\tfrac{1}{3}\rho\initscale^2v^2\tau^2\,(\mathrm{err}_0^2+\mathrm{err}_1^2).$ 
\end{lemma}

\begin{proof}
By \cref{lem:J-geometry}(a), $(\pd{\mathcal L}{c},\pd{\mathcal L}{g})^\top=\mathbf C_{\mathrm{grad}}(\mathrm{err}_0,\mathrm{err}_1)^\top$ with $\mathbf C_{\mathrm{grad}}=\bigl(\begin{smallmatrix}-1&1\\-r_0&r_1\end{smallmatrix}\bigr)$. Its singular values $s_1\ge s_2$
satisfy
 \[
\begin{gathered}
s_1s_2=|\det\mathbf C_{\mathrm{grad}}|=|\tau|,\\ s_1\le\|\mathbf C_{\mathrm{grad}}\|_F=(2+r_0^2+r_1^2)^{1/2}\le\sqrt3,\\ s_2\ge\frac{|\tau|}{\sqrt3} .
\end{gathered}
\] 
The second claim follows from \cref{lem:J-geometry}(d).
\end{proof}

\paragraph{Reflection.} The map $(c,g,s,a)\mapsto(-c,-g,s,-a)$ exchanges $m_0$ and $m_1$, and so maps $\mathcal L$ and \cref{eq:J-flow} for $p$ to those
for $1-p$. On parameters it is $(u_0,u_1,w,q,k)\mapsto(u_1,u_0,-w,-q,k)$, i.e.\ $q(0)\mapsto-q(0)$. It preserves $q(0)k(0)(1-2p)$, $\qkbalance$ and the
conclusion of \cref{thm:matched}. Hence we assume $p<\frac{1}{2}$, and set
 $h^\star=\log\frac{1-p}{p}>0,\qquad \mu=\frac{p(1-2p)}{14}.$ 
The initialization is \emph{aligned} if $q(0)k(0)<0$, i.e.\ $a(0)<0$.

\begin{lemma}[Drift identities]
\label{lem:J-drift}
\begin{enumerate}[label=(\alph*),leftmargin=*,itemsep=2pt,topsep=3pt]
\item $\operatorname{sign}\ddt{a}=-\operatorname{sign}g$.
\item $\ddt{Q}=-\frac{1}{2}\rho\initscale^2v^2\,\tau\,\mathrm{err}$, so $Q$ is nondecreasing while $a\le0$.
\item $(1+\frac{v^2}{4})\,g=Q+\frac{v^2}{2}\,h$.
\item $\ddt{h}=-\rho\initscale^2(1+\frac{v^2}{4})\pd{\mathcal L}{c}-\frac{1}{4}\rho\initscale^2v^2\,\tau\,\mathrm{err}$.
\item If $h\le\frac{h^\star}{2}$, then $\pd{\mathcal L}{c}\le0$. If $h\le\frac{h^\star}{4}$ and $|R|\le1$, then $\pd{\mathcal L}{c}\le-\mu$.
\end{enumerate}
\end{lemma}

\begin{proof}
(a) By \cref{eq:J-flow} and \cref{lem:J-geometry}(a), $\ddt{a}=-\frac{\initscale^4S}{d_{\mathrm h}}\,g\,\dd{\sigma(a)}{a}\,\mathrm{err}$ with $\dd{\sigma(a)}{a},\mathrm{err}>0$ and $S=q^2+k^2>0$ (if $\qkbalance=0$, then $S=\frac{2\sqrt{d_{\mathrm h}}|a|}{\initscale^2}$ and $a$ never reaches $0$ by \cref{rem:J-balance}).

(b), (d) By \cref{eq:J-flow} and \cref{lem:J-geometry}(a), $\ddt{Q}=-\rho\initscale^2v^2\bigl(\pd{\mathcal L}{g}-\frac{1}{2}\pd{\mathcal L}{c}\bigr)$ and
$\ddt{h}=\ddt{c}+\frac{1}{2}\ddt{g}$. (c) Substitute $c=h-\frac{g}{2}$.

(e) $m_1=m_0-2h$ and $\pd{\mathcal L}{c}=p\,\phi(m_1;s)-(1-p)\,\phi(m_0;s)$. By \cref{lem:J-geometry}(b),
\begin{alignat*}{2}
 \phi(m_1;s)&\le e^{2\max(h,0)}\phi(m_0;s) &\qquad& \why{$\phi$ nonincreasing; $\phi(m-2x)\le e^{2x}\phi(m)$}\\
 \pd{\mathcal L}{c}&\le-\phi(m_0;s)\bigl((1-p)-p\,e^{2\max(h,0)}\bigr),
\end{alignat*}
and $p\,e^{h^\star}=1-p$ gives the first claim. If $h\le\frac{h^\star}{4}$, then $p\,e^{2\max(h,0)}\le p\,e^{h^\star/2}=\sqrt{p(1-p)}$ and
\begin{alignat*}{2}
 (1-p)-\sqrt{p(1-p)}&=\frac{\sqrt{1-p}\,(1-2p)}{\sqrt{1-p}+\sqrt p}\ge\frac{1-2p}{2} &\qquad& \why{$\sqrt{1-p}\ge\tfrac{1}{\sqrt2}\ge\tfrac{1}{2}(\sqrt{1-p}+\sqrt p)$}\\
 \phi(m_0;s)&\ge\tfrac{1}{2}\sigma\bigl(-\tfrac{h^\star}{4}-\tfrac{1}{2}\bigr) && \why{$m_0=h+\tfrac{R}{2}\le\tfrac{h^\star}{4}+\tfrac{1}{2}$}\\
 &\ge\tfrac{1}{4}e^{-1/2}\Bigl(\frac{p}{1-p}\Bigr)^{1/4}\ge\frac{p}{7} && \why{$\sigma(-z)\ge\tfrac{1}{2}e^{-z}$}
\end{alignat*}
using $p^{1/4}\ge p$ and $\frac{1}{4}e^{-1/2}\ge\frac{1}{7}$. Multiply.
\end{proof}

\begin{lemma}[Alignment is kept]
\label{lem:J-align}
Let $t_0\ge0$ with
 $a(t_0)<0,\qquad g(t_0)\ge0,\qquad \bigl(1+\tfrac{v^2}{4}\bigr)g(t_0)\ge\tfrac{v^2}{2}\bigl(h(t_0)-\tfrac{h^\star}{2}\bigr)_+ .$ 
Then for all $t\ge t_0$: $g(t)\ge0$, $a(t)\le a(t_0)$, and $h(t)\ge h_0:=\min\{h(t_0),\frac{h^\star}{2}\}$.
\end{lemma}

\begin{proof}
Let $t'=\sup\{t\ge t_0:a<0\text{ on }[t_0,t]\}>t_0$. On $[t_0,t')$, $\tau<0$ and $\mathrm{err}>0$.

\emph{$h\ge h_0$.} Where $h\le\frac{h^\star}{2}$, $\ddt{h}>0$ by \cref{lem:J-drift}(d), (e). If $h(t_1)<h_0$ for some $t_1\in(t_0,t')$, let
$t_2=\sup\{t\le t_1:h(t)\ge h_0\}$. Then $h<h_0\le\frac{h^\star}{2}$ on $(t_2,t_1]$, so $h(t_1)>h(t_2)=h_0$, a contradiction.

\emph{$g\ge0$.} By \cref{lem:J-drift}(b), (c),
\begin{alignat*}{2}
 \bigl(1+\tfrac{v^2}{4}\bigr)g(t)&=\bigl(1+\tfrac{v^2}{4}\bigr)g(t_0)+\bigl(Q(t)-Q(t_0)\bigr) &\qquad&\\
 &\qquad+\tfrac{v^2}{2}\bigl(h(t)-h(t_0)\bigr)\\
 &\ge\bigl(1+\tfrac{v^2}{4}\bigr)g(t_0)-\tfrac{v^2}{2}\bigl(h(t_0)-h_0\bigr) && \why{$Q\uparrow$, $h\ge h_0$}\\
 &\ge0 && \why{$h(t_0)-h_0=(h(t_0)-\tfrac{h^\star}{2})_+$}
\end{alignat*}

\emph{$a\le a(t_0)$.} By \cref{lem:J-drift}(a), $\ddt{a}\le0$ on $[t_0,t')$. If $t'<\infty$, then $a(t')\le a(t_0)<0$ and, by continuity, $a<0$ on a
neighbourhood of $t'$, contradicting the definition of $t'$. So $t'=\infty$.
\end{proof}

\subsection{Conditions and scales}

\paragraph{Aligned threshold and scales.}\label{app:J-constants}
After reflection assume $p<1/2$. We use only $0<\varepsilon\le1/2$ and $q(0)k(0)<0$.
Let $a_{\mathrm{sc}}=\varepsilon^4/\sqrt{d_{\mathrm h}}$. The seed implies
\begin{equation}\label{eq:J-scales}
\begin{gathered}\tfrac14\varepsilon^2\le v^2\le4\varepsilon^2\le1,\quad
\tfrac14 a_{\mathrm{sc}}\le|a(0)|\le4a_{\mathrm{sc}}\le\tfrac14,\\
\sigma'(a)\ge\tfrac16\ (|a|\le1),\quad
\tanh(x/2)\ge x/4\ (0\le x\le2).\end{gathered}
\end{equation}
\subsection{Aligned gate growth and persistent margins}

The next two lemmas give a persistent readout floor and a time bound for $|a|$ to reach $1$.

Let $T_h=\frac{h^\star}{4\rho\initscale^2\mu}$ and $g_{\mathrm{low}}=\min\{1,\frac{v^2h^\star}{20}\}\ge\frac{\initscale^2}{80}\min\{1,h^\star\}$.

\begin{lemma}[Gate floor]
\label{lem:J-floor}
Let $t_0$ satisfy the hypotheses of \cref{lem:J-align}, with moreover $h(t_0)\ge0$ and $Q(t_0)\ge-\frac{v^2h^\star}{16}$. Then $g(t)\ge g_{\mathrm{low}}$ for all
$t\ge t_0+T_h$.
\end{lemma}

\begin{proof}
On $[t_0,\infty)$, $a<0$, $Q\ge Q(t_0)$ and $h\ge0$ (\cref{lem:J-align}, \cref{lem:J-drift}(b)). At any $s\ge t_0$ with $g(s)<g_{\mathrm{low}}$,
\begin{alignat*}{2}
 \tfrac{v^2}{2}h(s)&=\bigl(1+\tfrac{v^2}{4}\bigr)g(s)-Q(s)<\tfrac{5}{4}\cdot\tfrac{v^2h^\star}{20}+\tfrac{v^2h^\star}{16}=\tfrac{v^2h^\star}{8} &\qquad& \why{\cref{lem:J-drift}(c), $v^2\le1$}\\
 |R(s)|&=g(s)|\tau(s)|<g_{\mathrm{low}}\le1,
\end{alignat*}
so $h(s)<\frac{h^\star}{4}$, and by \cref{lem:J-drift}(d), (e) and $\tau(s)<0$,
 $\ddt{h}(s)\ge\rho\initscale^2\mu,\qquad \frac{d}{dt}\Bigl[\bigl(1+\tfrac{v^2}{4}\bigr)g\Bigr](s)=\ddt{Q}(s)+\tfrac{v^2}{2}\ddt{h}(s)>0 .$ 
Suppose $g(t)<g_{\mathrm{low}}$ for some $t\ge t_0+T_h$, and let $s_0=\sup\{s\in[t_0,t]:g(s)\ge g_{\mathrm{low}}\}$, or $s_0=t_0$ if there is none. On $(s_0,t]$, $g$
increases, so $g(t)>g(s_0)$; if $s_0>t_0$ then $g(s_0)=g_{\mathrm{low}}$, a contradiction. So $s_0=t_0$, and on $(t_0,t]$, $h$ increases at rate at least
$\rho\initscale^2\mu$: $h(t)\ge\rho\initscale^2\mu\,T_h=\frac{h^\star}{4}$, contradicting $h(t)<\frac{h^\star}{4}$.
\end{proof}

\begin{lemma}[Gate growth]
\label{lem:J-grow}
Let $t_0\ge0$ with $a(t_0)\le-\frac{a_{\mathrm{sc}}}{8}$, and suppose $g\ge g_{\mathrm{low}}$ on $[t_0,\infty)$. Then $t_3=\inf\{t\ge t_0:\mathcal L(t)\le\mathcal L^\star\}$ satisfies
 $t_3-t_0\le\frac{4080\sqrt{d_{\mathrm h}}\,\log\frac{8\sqrt{d_{\mathrm h}}}{\initscale^4}}{p_{\min}\min\{1,h^\star\}\,\initscale^4}+\frac{8010\,K\log K}{p_{\min}^2\initscale^4} .$ 
\end{lemma}

\begin{proof}
Since $g\ge0$, $a$ is nonincreasing (\cref{lem:J-drift}(a)), so $|a|$ is nondecreasing on $[t_0,\infty)$. Let $t_1=\inf\{t\ge t_0:|a(t)|\ge1\}$. On
$[t_0,\min\{t_1,t_3\})$,
\begin{alignat*}{2}
 \mathrm{err}&\ge\bigl(\mathrm{err}_0^2+\mathrm{err}_1^2\bigr)^{1/2}\ge\frac{p_{\min}}{17} &\qquad& \why{\cref{lem:J-levels}(c)}\\
 \frac{d|a|}{dt}&=\frac{\initscale^4S}{d_{\mathrm h}}\,g\,\dd{\sigma(a)}{a}\,\mathrm{err}\ge\frac{2\initscale^2|a|}{\sqrt{d_{\mathrm h}}}\cdot\frac{\initscale^2\min\{1,h^\star\}}{80}\cdot\frac{1}{6}\cdot\frac{p_{\min}}{17} && \why{$S\ge\frac{2\sqrt{d_{\mathrm h}}|a|}{\initscale^2}$, $|a|\le1$}\\
 &=\frac{\initscale^4\min\{1,h^\star\}\,p_{\min}}{4080\sqrt{d_{\mathrm h}}}\,|a|,
\end{alignat*}
so $\min\{t_1,t_3\}-t_0\le\frac{4080\sqrt{d_{\mathrm h}}}{\initscale^4\min\{1,h^\star\}p_{\min}}\log\frac{8}{a_{\mathrm{sc}}}$. If $t_1<t_3$, then on $[t_1,t_3)$, $|a|\ge1$ and
\begin{alignat*}{2}
 \ddt{\mathcal L}&\le-\tfrac{1}{3}\rho\initscale^2v^2\tanh^2(\tfrac{1}{2})\,(\mathrm{err}_0^2+\mathrm{err}_1^2) &\qquad& \why{\cref{lem:J-gradient}}\\
 &\le-\frac{\initscale^4p_{\min}^2}{8010\,K} && \why{\cref{lem:J-levels}(c), $\tanh^2\tfrac{1}{2}\ge\tfrac{1}{5}$}
\end{alignat*}
and $t_3-t_1\le\frac{8010K\log K}{p_{\min}^2\initscale^4}$, since $\mathcal L(t_1)\le\log K$.
\end{proof}

\begin{lemma}[Descent]
\label{lem:J-descent}
Suppose $|a(t)|\ge\frac{a_{\mathrm{sc}}}{8}$ for all $t\ge t_0$, and let $t_3=\inf\{t\ge t_0:\mathcal L(t)\le \mathcal L^\star\}$. Then
\[
 t_3-t_0\le\frac{1.7\cdot10^6\,K\log K\,d_{\mathrm h}}{p_{\min}^2\initscale^{12}},
\]
every prediction of $s_t$, $t\ge2$, has positive margin against all $K-1$ other counts for all $t\ge t_3$, and
$\mathcal L(t)\le\frac{2\cdot10^5Kd_{\mathrm h}}{\initscale^{12}(t-t_3)}$ for $t>t_3$.
\end{lemma}

\begin{proof}
For $t\ge t_0$,
\begin{alignat*}{2}
 \tau^2&\ge\tanh^2(\tfrac{a_{\mathrm{sc}}}{16})\ge\frac{a_{\mathrm{sc}}^2}{1024} &\qquad& \why{\cref{eq:J-scales}}\\
 \ddt{\mathcal L}&\le-\tfrac{1}{3}\rho\initscale^2v^2\tau^2\,(\mathrm{err}_0^2+\mathrm{err}_1^2) && \why{\cref{lem:J-gradient}}\\
 &\le-\frac{\initscale^{12}}{6144\,K\,d_{\mathrm h}}\,(\mathrm{err}_0^2+\mathrm{err}_1^2) && \why{$\rho=\tfrac{2}{K}$, $v^2\ge\tfrac{1}{4}\initscale^2$}
\end{alignat*}
While $\mathcal L>\mathcal L^\star$,
\begin{alignat*}{2}
 \ddt{\mathcal L}&\le-\frac{\initscale^{12}p_{\min}^2}{1640448\,K\,d_{\mathrm h}} &\qquad& \why{\cref{lem:J-levels}(c)}\\
 t_3-t_0&\le\frac{1640448\,K\log K\,d_{\mathrm h}}{p_{\min}^2\initscale^{12}} && \why{$\mathcal L(t_0)\le\mathcal L(0)=\log K$}
\end{alignat*}
For $t\ge t_3$, $\mathcal L<p_{\min}\log2$, so all margins are positive and $P_0,P_1>\frac{1}{2}$ (\cref{lem:J-levels}(a)), and
\begin{alignat*}{2}
 \ddt{\mathcal L}&\le-\frac{\initscale^{12}}{196608\,K\,d_{\mathrm h}}\,\mathcal L^2 &\qquad& \why{\cref{lem:J-levels}(b)}\\
 \frac{d}{dt}\frac{1}{\mathcal L}&\ge\frac{\initscale^{12}}{196608\,K\,d_{\mathrm h}} && \qedhere
\end{alignat*}
\end{proof}

\begin{proof}[Process part of \cref{thm:matched}]
\emph{Reduction.} By the reflection, $p<\frac{1}{2}$; then $p_{\min}=p$.

\emph{Aligned start.} $a(0)\le-\frac{a_{\mathrm{sc}}}{4}<0$ and $g(0)=h(0)=0$, so \cref{lem:J-align} applies at $t_0=0$: $a(t)\le-\frac{a_{\mathrm{sc}}}{4}$ for all $t$.
\Cref{lem:J-descent} with $t_0=0$ gives positive margins for $t\ge t_3$ and $\Lstep\to0$. \Cref{lem:J-floor} applies at $t_0=0$ ($h(0)=Q(0)=0$), and
\cref{lem:J-grow} with $t_0=T_h$ gives
 $t_3\le T_h+\frac{4080\sqrt{d_{\mathrm h}}\,\log\frac{8\sqrt{d_{\mathrm h}}}{\initscale^4}}{p\min\{1,h^\star\}\,\initscale^4}+\frac{8010\,K\log K}{p^2\initscale^4} .$ 


The explicit bound in \cref{eq:focused-time} is precisely this aligned estimate.
The loss is nonincreasing after crossing $\mathcal L^\star$; the strict threshold
in \cref{lem:J-levels}(a) therefore makes the margins persistent. The decay
estimate of \cref{lem:J-descent} proves $\Lstep\to0$.
\end{proof}
\paragraph{Scalar balance.}\label{rem:J-balance}
The identity $q^2-k^2=\Delta_{qk}$ is invariant. Even when it vanishes,
alignment persists: the equation for $a$ has the form $\dot a=-f(t)a$
on a finite time interval, with continuous $f$. Thus a nonzero aligned
$a(0)$ never reaches zero in finite time.
\paragraph{Readout restriction.}\label{rem:J-line}
The line $\mathbb R(I-R)$ is part of the trained family, not a conclusion
about unrestricted gate readouts. For $K=2$ it removes only a common-logit
component. For $K\ge3$, unrestricted circulant gate readouts introduce
additional background contributions to the gate derivative; this proof does
not cover them.


\section{Matching the initialization, metric, and time intervals}\label{app:match}
The supplied seed obeys \cref{ass:init} with $a_0=\max(4\sqrt K,\sqrt{2|\Sigma|})$, $c_{\mathrm{emb}}=1$ and $r_0=1/8$. In particular $\|\Pi_Q\|_{\mathrm{op}}=\sqrt K$, $\|\Pi_K\|_{\mathrm{op}}=\|\Pi_V\|_{\mathrm{op}}=1$, and all MLP-bias magnitudes equal $\varepsilon$. Zero embeddings for letters' format tokens and positions and a zero readout are allowed. A fixed causal submask changes only the input-independent probability weights at zero content scores. The estimates in Appendix~\ref{app:outcome} use their positivity and unit sum, so any such nonempty mask is covered.

There is a metric distinction that cannot be suppressed by calling the tied scalar family a block projection. Let $A$ embed its increments into ambient matrix parameters. The disjoint registers and
$\|\Pi_Q\|_F^2=K$, $\|\Pi_K\|_F^2=1$ give, for every tangent $\delta\vartheta$,
\begin{equation}\label{eq:metric-comparison}
\|\delta\vartheta\|^2\le\|A\delta\vartheta\|^2
=\|\delta U\|_F^2+\|\delta W\|_F^2+K(\delta q)^2+(\delta k)^2
\le K\|\delta\vartheta\|^2.
\end{equation}
For either objective, scalar Euclidean flow satisfies $\dot\vartheta=-A^\top\nabla_\theta\mathcal L$ and $\dot\theta=-AA^\top\nabla_\theta\mathcal L$. This is a positive semidefinite constant mobility, generally not an orthogonal projection. Its energy identity is exactly
\[
\frac{d\mathcal L}{dt}=-\|\dot\vartheta\|^2.
\]
The local loss floor is a statement about logits on the ambient ball, independent of the metric. Let $B_n=C_2\varepsilon^2\lambda^n$ and stop the outcome flow at the first ambient exit from radius $r_0\varepsilon$. Until that exit, the zero initial readout and \cref{eq:lo-floor} give
\[
\|\theta(t)-\theta_0\|^2\le K\|\vartheta(t)-\vartheta_0\|^2
\le Kt\int_0^t\|\dot\vartheta\|^2\,du\le KtB_n.
\]
Hence exit cannot occur before $r_0^2\lambda^{-n}/(KC_2)$. The factor $1/2$ in Theorem~\ref{thm:matched} provides an interior interval. The same loss identity yields $\E\KL(u\|p_t)\le B_n$, and therefore
\[
\big|\E[p_{t,y}]-1/K\big|\le\E\mathrm{TV}(p_t,u)
\le\sqrt{B_n/2}.
\]
This proves the stated outcome conclusions for exactly the same scalar metric as the process proof. The absolute-value probability form follows because each coordinate's absolute deviation from $1/K$ is bounded by total variation. The broader ambient orthogonal-block flow proved in Appendix~\ref{app:outcome} remains a separate supporting argument; it does not give unrestricted positive training.

For fixed $K,p$, $N_b=\Theta(\sqrt n)$ at the executor MLP width, so $\varepsilon=c_*n^{-1/2}$ satisfies \cref{eq:small} for sufficiently small fixed $c_*$ and sufficiently large $n$. The constants in the local loss floor are width-independent in this regime. In \cref{eq:focused-time}, $T_h=O(n)$, the gate-growth term is $O(n^{5/2}\log n)$, and the remaining readout term is $O(n^2)$. Thus $T_{\mathrm{step}}=O(n^{5/2}\log n)$ while $T_{\mathrm{out}}=c\lambda^{-n}$, with $c>0$ independent of $n$. Since an exponential dominates this polynomial, their intervals overlap for all sufficiently large $n$. At any finite horizon the exact inequality $T_{\mathrm{step}}<T_{\mathrm{out}}$ is required; merely quoting asymptotic orders does not establish overlap.

\section{Scope of the fixed query--key refinement}\label{app:robustness}
The authoritative 64-page draft additionally proves a fixed-perturbation refinement in its Appendix D.3. Its assumptions are $K=2$, $0<p<1/2$, aligned $q_0k_0<0$, $0<\varepsilon\le p$, the same supplied mask and fixed value map, and
\begin{equation}\label{eq:planted-perturbation}
\begin{gathered} W_Q=q\Pi_Q+\Delta_Q,\quad W_K=k\Pi_K+\Delta_K,\\
\langle\Delta_K,\Pi_K\rangle_F=0,\qquad
\|\Delta_Q\|_{\mathrm{op}}\le|q_0|/8,\quad
\|\Delta_K\|_{\mathrm{op}}\le|k_0|/8.\end{gathered}
\end{equation}
Only $(U,W,q,k)$ train; $\Delta_Q,\Delta_K$ remain fixed. The reported positive bound is
\[
T\le\frac C{\varepsilon^4}\left(
\frac{\sqrt{d_{\mathrm h}}\log(d_{\mathrm h}/\varepsilon)}{p^2\min\{1,h^\star\}}
+\frac1{p^4}\right),
\]
with persistent margins and loss tending to zero. This refinement uses the same training family with fixed perturbed maps. It is not needed for Theorem~\ref{thm:matched}; its complete source proof follows. It does not license arbitrary perturbations, trained perturbations, value perturbations, general $K$, or unaligned starts. \paragraph{The perturbed gate.} Without a value perturbation the readout features are those of \cref{lem:J-features};
only the gate changes.

\begin{lemma}[Perturbed gate]
\label{lem:perturbed-gate}
Under \cref{eq:planted-perturbation}, let $\bm\delta_s=\bm\Delta_Q\hat{\bm e}^{\mathrm c}_s$, $\beta_s=\langle\hat{\bm e}_{\mathrm h},\bm\delta_s\rangle$,
$\gamma_s=\langle\bm\delta_s,\bm\Delta_K\hat{\bm e}_\ell\rangle$, $\alpha_s=\langle\hat{\bm e}_{\mathrm h},\bm\Delta_K(\hat{\bm e}_0-\hat{\bm e}^{\mathrm c}_s)\rangle$,
$\omega_s=\langle\bm\delta_s,\bm\Delta_K(\hat{\bm e}_0-\hat{\bm e}^{\mathrm c}_s)\rangle$ and $\kappa=\initscale^2/\sqrt{d_{\mathrm h}}$. The prediction of $s_t$ with
$s_{t-1}=s$ and $b_t=b$ has logits $\initscale\mathbf U\bm e_s+\initscale v\,r^s_b\,\mathbf W\bm e_s$, where
\begin{equation}
 r^s_b=\sigma\bigl(-(-1)^ba_s-o_s\bigr),\qquad a_s=\kappa\bigl(k(q+\beta_s)+\gamma_s\bigr),\qquad o_s=\kappa(q\alpha_s+\omega_s),
 \label{eq:perturbed-gate}
\end{equation}
and
\begin{equation}
 r^s_0-r^s_1=\frac{\sinh(-a_s)}{\cosh a_s+\cosh o_s}.
 \label{eq:perturbed-contrast}
\end{equation}
With operator norms, $|\beta_s|\le\|\bm\Delta_Q\|$, $|\gamma_s|\le\|\bm\Delta_Q\|\|\bm\Delta_K\|$, $|\alpha_s|\le\sqrt2\|\bm\Delta_K\|$ and
$|\omega_s|\le\sqrt2\|\bm\Delta_Q\|\|\bm\Delta_K\|$.
\end{lemma}

\begin{proof}
As in \cref{lem:J-features}, block $1$ and the MLPs add zero. Since $\bm\Pi_Q\hat{\bm e}^{\mathrm c}_s=\hat{\bm e}_{\mathrm h}$,
$\bm\Pi_K\hat{\bm e}^{\mathrm c}_s=0$ and $\bm\Pi_K(\hat{\bm e}_0+(-1)^b\hat{\bm e}_\ell)=(-1)^b\hat{\bm e}_{\mathrm h}$, the query at $i_t$ and the
two keys are
\[
 \initscale(q\hat{\bm e}_{\mathrm h}+\bm\delta_s),\qquad \mathrm{key}_{i_t}=\initscale\bm\Delta_K\hat{\bm e}^{\mathrm c}_s,\qquad
 \mathrm{key}_{j_t}=\initscale\bigl((-1)^b(k\hat{\bm e}_{\mathrm h}+\bm\Delta_K\hat{\bm e}_\ell)+\bm\Delta_K\hat{\bm e}_0\bigr).
\]
The score difference $\mathrm{score}(j_t)-\mathrm{score}(i_t)$ is $\frac{\initscale^2}{\sqrt{d_{\mathrm h}}}$ times the inner product of the query with
$\mathrm{key}_{j_t}-\mathrm{key}_{i_t}$, which is $(-1)^ba_s+o_s$ because $\langle\hat{\bm e}_{\mathrm h},\bm\Delta_K\hat{\bm e}_\ell\rangle=\langle\bm\Delta_K,\bm\Pi_K\rangle_F=0$.
So $i_t$ gets weight $r^s_b$. The values are $\initscale v\hat{\bm e}^{\mathrm g}_s$ at $i_t$ and $0$ at $j_t$, so the residual and logits are those of
\cref{lem:J-features} with $r_b$ replaced by $r^s_b$. For \cref{eq:perturbed-contrast}, put $x=-a_s$, $o=o_s$; then
$\sigma(x-o)-\sigma(-x-o)=\frac{e^{x-o}-e^{-x-o}}{(1+e^{x-o})(1+e^{-x-o})}=\frac{2e^{-o}\sinh x}{2e^{-o}(\cosh x+\cosh o)}$. The bounds use
$\|\bm\delta_s\|\le\|\bm\Delta_Q\|$ and $\|\hat{\bm e}_0-\hat{\bm e}^{\mathrm c}_s\|=\sqrt2$.
\end{proof}

By \cref{eq:perturbed-contrast}, a state's letter contrast is positive exactly when $a_s<0$, whatever the offset; when $|o_s|\le|a_s|$ it is at
least $\frac12\tanh|a_s|$. The following source refinement controls $a_s$ and $o_s$ state by state, relative to the planted strength
\[
 X:=\kappa\,q|k|.
\]

\paragraph{Setting of the fixed-perturbation guarantee.} $K=2$, $0<p<\frac12$, $h^\star=\log\frac{1-p}{p}$, an aligned start $q(0)k(0)<0$ with
$\frac{\initscale}{2}\le|q(0)|,|k(0)|,|v|\le2\initscale$, zero readouts, and
\begin{equation}
 \|\bm\Delta_Q\|\le\tfrac18|q(0)|,\qquad \|\bm\Delta_K\|\le\tfrac18|k(0)|,\qquad 0<\initscale\le p .
 \label{eq:perturbed-radius}
\end{equation}
Since $\|\bm\Pi_Q\|_F=\sqrt2$ and $\|\bm\Pi_K\|_F=1$, the Frobenius condition $\|\bm\Delta_X\|_F\le\rho\,|x_0|\,\|\bm\Pi_X\|_F$ implies
\cref{eq:perturbed-radius} for $\rho\le\frac{1}{8\sqrt2}$. Changing the signs of $\mathbf W_Q$ and $\mathbf W_K$ together, perturbations included, leaves every score
and \cref{eq:perturbed-radius} unchanged, so we take $q(0)>0>k(0)$; write $\delta=\|\bm\Delta_Q\|$ and $\eta_K=\|\bm\Delta_K\|$.

\begin{lemma}[Reduced flow for $K=2$]
\label{lem:perturbed-reduce}
Let $c_s=\initscale(U_{ss}-U_{s+1,s})$ and $g=2\initscale vw$. The step loss is
$L=\frac12\sum_{s\in\{0,1\}}L_s$ with $L_s=(1-p)\,\ell(m_{s0})+p\,\ell(m_{s1})$, $\ell(m)=\log(1+e^{-m})$, and margins
$m_{s0}=c_s+g\,r^s_0$, $m_{s1}=-c_s-g\,r^s_1$. With $E_{s0}=(1-p)\sigma(-m_{s0})$, $E_{s1}=p\,\sigma(-m_{s1})$,
$\Lambda_s=E_{s0}r^s_0(1-r^s_0)+E_{s1}r^s_1(1-r^s_1)$ and $\Gamma_s=E_{s0}r^s_0(1-r^s_0)-E_{s1}r^s_1(1-r^s_1)$, the flow is
\begin{gather*}
 \ddt{c_s}=\initscale^2(E_{s0}-E_{s1}),\qquad \ddt{g}=\frac{\initscale^2v^2}{2}\sum_s\bigl(E_{s0}r^s_0-E_{s1}r^s_1\bigr),\\
 \ddt{q}=-\frac{\kappa g}{2}\sum_s\bigl(\Lambda_sk+\Gamma_s\alpha_s\bigr),\qquad \ddt{k}=-\frac{\kappa g}{2}\sum_s\Lambda_s(q+\beta_s),
\end{gather*}
and the other two coordinates of $\mathbf U$ receive no gradient.
\end{lemma}

\begin{proof}
For $K=2$ the cross-entropy depends only on the difference of the two logits, which by \cref{lem:perturbed-gate} is $c_s+g\,r^s_b$ for the
prediction from state $s$; the target is $s+b$. The pair $(s_{t-1},b_t)$ is uniform $\times$ Bernoulli$(p)$ for every $t\ge2$, which gives $L$.
The coordinates $U_{ss}+U_{s+1,s}$ do not enter $L$. On $\mathbf U$ the Frobenius gradient gives $\ddt{c_s}=-2\initscale^2\pd{L}{c_s}$, and
$\|w(\mathbf I-\mathbf R)\|_F^2=4w^2$ gives $\ddt{g}=-\initscale^2v^2\pd{L}{g}$. Differentiating, $\pd{L_s}{c_s}=E_{s1}-E_{s0}$,
$\pd{L_s}{g}=E_{s1}r^s_1-E_{s0}r^s_0$, $\pd{L_s}{a_s}=g\Lambda_s$ and $\pd{L_s}{o_s}=g\Gamma_s$; with $\pd{a_s}{q}=\kappa k$, $\pd{a_s}{k}=\kappa(q+\beta_s)$,
$\pd{o_s}{q}=\kappa\alpha_s$ and $\pd{o_s}{k}=0$ this gives the gate equations.
\end{proof}

Note that $q^2-k^2$ is not conserved and $r^s_0+r^s_1\ne1$ when $o_s\ne0$; the next three lemmas replace the uses of these facts in
\cref{lem:J-align,lem:J-floor}.

\begin{lemma}[Statewise invariant region]
\label{lem:perturbed-invariant}
For all $t>0$: $q$ is nondecreasing and $k$ nonincreasing, so $q\ge q(0)$, $|k|\ge|k(0)|$ and $X$ is nondecreasing; $g>0$; and $c_s<h^\star$
for both $s$. Moreover, for both $s$,
\begin{equation}
 \begin{gathered}
 \tfrac{55}{64}X\le-a_s\le\tfrac{73}{64}X,\qquad |o_s|\le\tfrac{9\sqrt2}{64}X\le\tfrac15X,\\
 r^s_0\ge\tfrac12\ge r^s_1,\qquad r^s_0-r^s_1\ge\tfrac12\tanh|a_s|>0.
 \end{gathered}
 \label{eq:perturbed-gate-bounds}
\end{equation}
The flow exists for all $t\ge0$.
\end{lemma}

\begin{proof}
\emph{Gate bounds.} Wherever $q\ge q(0)\ge8\delta$ and $|k|\ge|k(0)|\ge8\eta_K$, \cref{lem:perturbed-gate} gives
$-a_s=\kappa(|k|(q+\beta_s)-\gamma_s)\in\kappa q|k|\,[1-\frac18-\frac1{64},\,1+\frac18+\frac1{64}]$ and
$|o_s|\le\kappa\sqrt2\eta_K(q+\delta)\le\kappa\sqrt2\frac{|k|}8\cdot\frac{9q}8$, which is \cref{eq:perturbed-gate-bounds} for $a_s,o_s$. Then $|o_s|<|a_s|$,
so $r^s_0=\sigma(|a_s|-o_s)\ge\frac12\ge\sigma(-|a_s|-o_s)=r^s_1$, and \cref{eq:perturbed-contrast} with $\cosh o_s\le\cosh a_s$ gives the contrast bound.

\emph{Gate monotonicity on intervals with $g\ge0$.} Since $|\Gamma_s|\le\Lambda_s$, $|\alpha_s|\le\sqrt2\eta_K$ and $k<0$,
\[
 \begin{aligned}
 \ddt{q}&=\frac{\kappa g}{2}\sum_s\bigl(\Lambda_s|k|-\Gamma_s\alpha_s\bigr)
 \ge\frac{\kappa g}{2}\bigl(|k|-\sqrt2\eta_K\bigr)\sum_s\Lambda_s,\\
 \ddt{|k|}&=\frac{\kappa g}{2}\sum_s\Lambda_s(q+\beta_s)
 \ge\frac{\kappa g}{2}(q-\delta)\sum_s\Lambda_s.
 \end{aligned}
\]
Both are nonnegative while $|k|>\sqrt2\eta_K$ and $q>\delta$. These strict conditions hold at $t=0$, and by the displayed signs they cannot first fail
inside an interval on which $g\ge0$; so on any such interval $q\ge q(0)$, $|k|\ge|k(0)|$, both are monotone, and the gate bounds hold.

\emph{Readout region.} At $t=0$, $g=0$, $c_s=0$ and $E_{s0}=\frac{1-p}2>\frac p2=E_{s1}$, so $\ddt{g}(0)>0$ by $r^s_0>r^s_1$; hence
$t^\ast=\inf\{t>0:g(t)\le0\text{ or }\max_sc_s(t)\ge h^\star\}>0$. Suppose $t^\ast<\infty$; on $[0,t^\ast]$, $g\ge0$ and the gate bounds hold.
If $g(t^\ast)=0$, then $\ddt{g}(t^\ast)\le0$. But at $g=0$, $E_{s0}r^s_0-E_{s1}r^s_1=(1-p)\sigma(-c_s)r^s_0-p\,\sigma(c_s)r^s_1>0$, because
$p\,\sigma(c_s)\le(1-p)\sigma(-c_s)$ for $c_s\le h^\star$ and $r^s_0>r^s_1$; a contradiction. Otherwise $g(t^\ast)>0$ and $c_s(t^\ast)=h^\star$ for some
$s$, so $\ddt{c_s}(t^\ast)\ge0$. But $\sigma(-h^\star)=p$, so $E_{s0}=(1-p)\sigma(-h^\star-gr^s_0)<p(1-p)\le p\,\sigma(h^\star+gr^s_1)=E_{s1}$, a contradiction.
So $t^\ast=\infty$.

\emph{Existence.} $|\ddt{c_s}|\le\initscale^2$ and $|\ddt{g}|\le\initscale^2v^2$, so $c_s,g$ grow at most linearly; given $g$, the $(q,k)$ equations are
linear-growth, so Gr\"onwall excludes blow-up in finite time.
\end{proof}

\begin{lemma}[Balance and readout floor]
\label{lem:perturbed-floor}
Let $Q=g-\frac{v^2}4(c_0+c_1)$, $\mu=\frac{p(1-2p)}3$, $g_{\mathrm{low}}=\frac{v^2h^\star}{16}$ and $T_h=\frac{h^\star}{2\initscale^2\mu}$.
On every interval $[0,t]$ on which $X\le1$, $Q$ is nondecreasing and $Q\ge0$, and $g(t)\ge g_{\mathrm{low}}$ if $t\ge T_h$.
\end{lemma}

\begin{proof}
\emph{Balance.} Write $E_{s0}r^s_0-E_{s1}r^s_1=\bar r_s(E_{s0}-E_{s1})+\frac12(r^s_0-r^s_1)(E_{s0}+E_{s1})$ with $\bar r_s=\frac12(r^s_0+r^s_1)$.
Since $|\partial_o(\sigma(-a-o)+\sigma(a-o))|\le\frac12$, $|\bar r_s-\frac12|\le\frac{|o_s|}4$, and so
\[
 \begin{aligned}
 \ddt{Q}&=\frac{\initscale^2v^2}2\sum_s\Bigl[
 \bigl(\bar r_s-\tfrac12\bigr)(E_{s0}-E_{s1})
 +\tfrac12(r^s_0-r^s_1)(E_{s0}+E_{s1})\Bigr]\\
 &\ge\frac{\initscale^2v^2}4\sum_s(E_{s0}+E_{s1})
 \Bigl(r^s_0-r^s_1-\frac{|o_s|}2\Bigr).
 \end{aligned}
\]
For $X\le1$, \cref{eq:perturbed-gate-bounds} gives $|a_s|\le\frac{73}{64}$, so $r^s_0-r^s_1\ge\frac12\tanh|a_s|\ge0.35\,|a_s|\ge0.3X\ge\frac{|o_s|}2$.
Hence $\ddt{Q}\ge0$, and $Q(0)=0$.

\emph{Floor.} For $c\le h^\star$, $0\le g\le1$ and $r_0,r_1\in[0,1]$, $E_0-E_1\ge f(c+g)$ with $f(y)=(1-p)\sigma(-y)-p\,\sigma(y)$, which is
decreasing, has $|\dd{f}{y}|\le\frac14$, $f(h^\star)=0$, and
$f(\frac{h^\star}2)=\frac{\sqrt{p(1-p)}(1-2p)}{1+2\sqrt{p(1-p)}}\ge\mu$. So $\ddt{c_s}\ge\initscale^2\mu$ when $c_s+g\le\frac{h^\star}2$, and
$\ddt{c_s}\ge-\frac{\initscale^2g}4$ always. By \cref{eq:perturbed-radius} and $h^\star\le\frac{1-2p}{p}$,
$g_{\mathrm{low}}\le\frac{\initscale^2h^\star}4\le\min\{1,\frac{3h^\star}8,2\mu\}$.

Suppose $t\ge T_h$, $X\le1$ on $[0,t]$ and $g(t)<g_{\mathrm{low}}$. Let $s_0=\sup\{s'\le t:g(s')\ge g_{\mathrm{low}}\}$, or $s_0=0$ if there is none.
On $(s_0,t]$, $g<g_{\mathrm{low}}$ and $Q\ge0$ give $c_0+c_1\le\frac{4g}{v^2}<\frac{h^\star}4$. Some state then has $c_s<\frac{h^\star}8$ and
$c_s+g<\frac{h^\star}2$, so with the other state's lower bound
\[
 \ddt{}(c_0+c_1)\ge\initscale^2\bigl(\mu-\tfrac{g_{\mathrm{low}}}4\bigr)\ge\tfrac12\initscale^2\mu,\qquad
 \ddt{g}=\ddt{Q}+\tfrac{v^2}4\ddt{}(c_0+c_1)>0 .
\]
If $s_0>0$, then $g(s_0)=g_{\mathrm{low}}$ and $g$ increases on $(s_0,t]$, contradicting $g(t)<g_{\mathrm{low}}$. If $s_0=0$, then
$c_0+c_1\ge\frac12\initscale^2\mu\,T_h=\frac{h^\star}4$ at $t$, a contradiction.
\end{proof}

\begin{lemma}[Error floor and gate growth]
\label{lem:perturbed-grow}
Let $z=\frac p4\log2$ and $e_\ast=\frac{p^2}{12}$. If $L>z$, then $\sum_{s,b}E_{sb}\ge e_\ast$. If moreover $t\ge T_h$ and $X\le1$, then
$\ddt{X}\ge\frac{\kappa\,g_{\mathrm{low}}\,p^2}{105}\,X$.
\end{lemma}

\begin{proof}
Each of the four (state, letter) cases has weight $w_{sb}=\frac{\pi_b}2\in[\frac p2,\frac12]$. If $L>z$, some case has $w_{sb}\ell(m_{sb})\ge\frac L4$,
so $\ell(m_{sb})>\frac z2$; since $\sigma(-m)=1-e^{-\ell(m)}$ and $\frac z2<0.05$, $E_{sb}\ge p(1-e^{-z/2})\ge0.97\,p\frac z2\ge e_\ast$.

For the growth, \cref{lem:perturbed-invariant} gives $|k|-\sqrt2\eta_K\ge(1-\frac{\sqrt2}8)|k|$ and $q-\delta\ge\frac78q$, so
\[
 \ddt{(q|k|)}=|k|\ddt{q}+q\ddt{|k|}\ge\frac{\kappa g}2\cdot\frac45(q^2+k^2)\sum_s\Lambda_s\ge\frac45\kappa g\,q|k|\sum_s\Lambda_s .
\]
For $X\le1$ every attention argument satisfies $|a_s|+|o_s|\le\frac{73+9\sqrt2}{64}<1.34$, so $r(1-r)\ge\dd{\sigma(x)}{x}\big|_{x=1.34}\ge\frac17$ and
$\sum_s\Lambda_s\ge\frac17\sum_{s,b}E_{sb}\ge\frac{e_\ast}7$. With $g\ge g_{\mathrm{low}}$ from \cref{lem:perturbed-floor},
$\ddt{X}\ge\frac45\cdot\frac17\cdot\frac{p^2}{12}\kappa g_{\mathrm{low}}X$.
\end{proof}

\begin{lemma}[Statewise descent]
\label{lem:perturbed-descent}
If $X\ge1$, then $\ddt{L}\le-\frac{\initscale^4}{6732}\bigl(\sum_{s,b}E_{sb}\bigr)^2$.
\end{lemma}

\begin{proof}
The readout part of the flow alone dissipates $2\initscale^2\sum_s(\pd{L}{c_s})^2+\initscale^2v^2(\pd{L}{g})^2
=\frac{\initscale^2}2D^2+\frac{\initscale^2v^2}4G^2$, with $D^2=\sum_s(E_{s0}-E_{s1})^2$ and $G=\sum_s(E_{s0}r^s_0-E_{s1}r^s_1)$; the gate part only adds
dissipation. For $X\ge1$, \cref{eq:perturbed-gate-bounds} gives $r^s_0-r^s_1\ge b:=\frac12\tanh\frac{55}{64}\ge\frac13$. Write $S_0=\sum_sE_{s0}$.
Since $E_{s0}r^s_0-E_{s1}r^s_1=E_{s0}(r^s_0-r^s_1)+r^s_1(E_{s0}-E_{s1})$, $G\ge bS_0-\sqrt2D$. If $D\ge\frac{bS_0}{2\sqrt2}$ the first term is at least
$\frac{\initscale^2b^2S_0^2}{16}$; otherwise $G\ge\frac{bS_0}2$ and the second is at least $\frac{\initscale^2v^2b^2S_0^2}{16}$. As $v^2\le1$, the dissipation is at
least $\frac{\initscale^2v^2b^2}{16}\max\{S_0,D\}^2$, and $\sum_{s,b}E_{sb}\le2S_0+\sqrt2D\le(2+\sqrt2)\max\{S_0,D\}$. Use $v^2\ge\frac{\initscale^2}4$ and
$b^2\ge\frac19$.
\end{proof}

\begin{proof}[Proof of the fixed binary refinement]
By \cref{lem:perturbed-invariant} the flow exists, stays in the aligned region and $X$ is nondecreasing, with
$X(0)\ge\frac{\initscale^4}{4\sqrt{d_{\mathrm h}}}$ and $X(0)\le\frac{4\initscale^4}{\sqrt{d_{\mathrm h}}}<1$. Let $t_3=\inf\{t:L(t)\le z\}$.

\emph{Growth.} Let $t_1=\inf\{t:X(t)\ge1\}$. On $[T_h,\min\{t_1,t_3\})$, \cref{lem:perturbed-grow} gives exponential growth of $X$ at rate
$\frac{\kappa g_{\mathrm{low}}p^2}{105}\ge\frac{\initscale^4h^\star p^2}{6720\sqrt{d_{\mathrm h}}}$, so
$\min\{t_1,t_3\}\le T_h+\frac{6720\sqrt{d_{\mathrm h}}}{\initscale^4h^\star p^2}\log\frac{4\sqrt{d_{\mathrm h}}}{\initscale^4}$.

\emph{Descent.} If $t_1<t_3$, then $X\ge1$ on $[t_1,\infty)$, and \cref{lem:perturbed-descent,lem:perturbed-grow} give
$\ddt{L}\le-\frac{\initscale^4e_\ast^2}{6732}$ while $L>z$; since $L(t_1)\le\log2$, $t_3-t_1\le\frac{6732\cdot144\log2}{\initscale^4p^4}$.

\emph{Margins.} If $L\le z$, every case has $\ell(m_{sb})\le\frac{z}{w_{sb}}\le\frac{\log2}2<\log2$, so every margin is positive; $L$ never increases,
so this persists for all $t\ge t_3$. After $t_3$, all margins are positive, so
$\ell(m)\le2\sigma(-m)$ and $L\le\sum_{s,b}E_{sb}$.
If $t_1<\infty$, put $t_4=\max\{t_1,t_3\}$. For $t\ge t_4$,
\cref{lem:perturbed-descent} applies and gives
\[
 \ddt{L}\le-\frac{\initscale^4}{6732}L^2,\qquad
 L(t)\le\frac{6732}{\initscale^4(t-t_4)}\longrightarrow0
 \quad(t>t_4).
\]
Otherwise $X<1$ for all time. If the nonincreasing loss had limit
$L_\infty>0$, then for $t\ge\max\{T_h,t_3\}$,
\cref{lem:perturbed-floor} and the growth calculation in
\cref{lem:perturbed-grow} would give
\[
 \sum_s\Lambda_s\ge\frac17\sum_{s,b}E_{sb}\ge\frac{L_\infty}{7},\qquad
 \ddt{X}\ge\frac{4\kappa g_{\mathrm{low}}L_\infty}{35}X.
\]
Exponential growth contradicts $X<1$. Hence $L_\infty=0$ in this case too.

\emph{Time.} $T_h=\frac{3h^\star}{2\initscale^2p(1-2p)}\le\frac{3}{2\initscale^2p^2}$, so
\[
 t_3\le\frac{3}{2\initscale^2p^2}+\frac{6720\sqrt{d_{\mathrm h}}}{\initscale^4h^\star p^2}\log\frac{4\sqrt{d_{\mathrm h}}}{\initscale^4}+\frac{7\cdot10^5}{\initscale^4p^4}
 \le\frac{C}{\initscale^4}\Bigl(\frac{\sqrt{d_{\mathrm h}}\log\frac{d_{\mathrm h}}{\initscale}}{p^2\min\{1,h^\star\}}+\frac1{p^4}\Bigr).\qedhere
\]
\end{proof}

The radius in \cref{eq:perturbed-radius} is a constant: it does not shrink with $\initscale$, $d_{\mathrm h}$ or $n$. The fixed perturbation is also
relatively suppressed during training, since $q$ and $|k|$ only grow.



\section{Experimental protocols and evidence}\label{app:experiments}
\subsection{Controlled counting}
The displacement measurements use two blocks of residual width $16$, three heads of dimension $8$, and MLP width $16$. Binary letters have match probability $1/4$; starts are uniform. At $\varepsilon=0.002$ the scaled bias coordinates have magnitude at least one and the readout is zero. Gradients enumerate all words through $n=14$ in float64. Seeds $0,1,2$ change the scaled parameter draw. The step objective averages transitions $t=2,\ldots,n$.

Finite updates are normalized to prescribed displacement and their step sizes accumulated as approximate flow time. The larger diagnostic radius $0.1\varepsilon\|\tilde\theta\|$ exceeds $r_0\varepsilon$: the scaled norm is about $70$, giving $r_0/\|\tilde\theta\|\simeq1.8\times10^{-3}$. Larger-radius runs use increments $10^{-3}\varepsilon\|\tilde\theta\|$; source checks halving or quartering them changed times by less than $1\%$. The theorem-radius plot uses stored $10^{-6}$ increments; $10^{-5}$ changed times by less than $0.5\%$. The logs record no ReLU sign changes before either exit. The figure reports medians and full seed ranges, not confidence intervals. The exponential guide is anchored at $n=4$; its vertical scale is not a numerical theorem constant. Compact widths test the plateau mechanism and need not contain the exact executors.

The $n=16$ AdamW comparison uses the executor architecture with all weights trainable, four letters, $20{,}000$ training and $1{,}000$ test prompts, Gaussian tensors of standard deviation $0.02$, and zero readout. Learning rate is $10^{-3}$, clipping $1$, batch $256$, budget $8$k updates, and paired seeds $42$--$51$. Step supervision includes all continuation and formatting labels. These runs change ReLU patterns and lie outside the theorem.

\subsection{Finite-step routing}
The source's $17$ pairs use counting modulo $2$, $p=1/4$, fan-in initialization at $\varepsilon=0.3$ and full-batch gradient descent with step size $2$. Query--key matrices, relative biases and readout train; values, embeddings and MLPs remain fixed. Each pair shares a seed and a $300$k update limit. Seeds are $0$--$4$ at $n=6,8,12$ and $0,1$ at $n=16$. Solved means every held-out prediction is correct. The process objective includes $t=1$, outside the main theorem's loss. At $n=6,8,12,16$, process solves $5/5,5/5,4/5,2/2$; all outcome pairs remain unsolved. Step median solve times are $44,61,77,237$ thousand updates. These source summaries are empirical evidence, not a newly rerun sweep or a verification of the continuous-time bound.

\subsection{Boolean-8 Qwen LoRA}
The task starts from four bits and applies eight randomly wired NOT, CNOT, SWAP or Toffoli gates. The pretrained model and tokenizer differ from the two-block architecture. \texttt{src/lora.py} trains rank-$8$ LoRA, $\alpha=16$, zero dropout, no bias adaptation, and target modules \texttt{all-linear}; base weights are frozen. The recorded model has $4{,}399{,}104$ trainable adapter parameters. AdamW uses learning rate $3\times10^{-4}$, zero weight decay, clipping $1$, batch $16$, and $800$ updates. The corpus has $12{,}000$ training and $300$ validation prompts; maximum sequence length is $384$. Seeds $2001$--$2003$ change initialization, not corpora.

Outcome predicts the final answer directly; process displays running states before it; answer-first places the answer before the trace. The eightfold outcome control uses $6400$ updates and the same batch size. It does not match tokens or prefix information. Table~\ref{tab:transfer} is extracted from the individual one-row CSVs under \texttt{results/paper/qwen05/main/} and \texttt{outcome\_6400/}. SD uses denominator $3-1$. All these runs lie outside the theorem's model, seed and optimizer.

\subsection{BIG-Bench Hard object tracking}
For $m=3,5,7$, \texttt{src/shuffled.py} generates $12{,}000$ training items from benchmark scenes with fresh assignments and swaps. The generated test set contains $300$ items. Separate official testing uses $250$ unchanged BIG-Bench Hard items per size from \texttt{data/bbh/}. The parser reconstructs those items exactly. We extract nested \texttt{bbh.answer\_accuracy}, not top-level generated-test accuracy.

Qwen adapters use rank $8$, $800$ updates, batch $16$, learning rate $3\times10^{-4}$, zero weight decay and clipping $1$. Traces list each holder's item and option letter after every swap. Item-only traces in earlier diagnostics stayed near the option-letter base rate and are not mixed into these results. Zero-shot official accuracies are $33.6,16.4,11.2\%$. Three-object outcome runs choose a common option with base rate $35.2\%$. The failing five-object process seed is $2003$: $46\%$ on official items and about $45\%$ on generated items. Including it gives process mean $82.0\%$, below outcome mean $86.0\%$. The manuscript retains all sizes and seeds. These pretrained-model experiments provide qualitative transfer evidence, with no claim that counting's bound applies to object tracking.

\subsection{Sparse count-label control}\label{app:label-budget}
\texttt{Paper/research/label\_budget/run\_control.py} uses count CE plus colon CE plus EOS CE, divided by three. Full process averages all count labels (running states and final copy); sparse process uniformly selects one count label per prompt/update, with an independent label RNG; outcome has one final-count label. Initialization, prompts, minibatch order and optimizer agree. Full and sparse process gradients agree in expectation on their teacher-forced rows; label sampling changes variance. The objective reweights count/format labels relative to the primary token-averaged experiments.

The reported sweep uses $n=6,8,12$, seeds $42$--$51$, $30$k updates, Gaussian weights of standard deviation $0.02$, zero readout, AdamW learning rate $0.001$, zero weight decay, batch $256$, clipping $1$, $20$k training and $1$k test prompts. Local evidence preserves all per-seed summaries. Sparse-process means are $97.14,99.83,99.19\%$, with sample SD $9.04,0.54,2.56$ points. Outcome means are $68.5,64.32,51.22\%$. Sparse process solves $9/10$ at each horizon at the $99\%$ threshold; outcome solves none. Full-process means are $97.0,89.1,100.0\%$, with $24/30$ total solved runs.

The initial $8$k, three-seed $n=8$ sweep was inconclusive: full process $70.5\%$, sparse process $62.1\%$, outcome $48.8\%$. It is retained in local evidence. The longer sweep supports an advantage under equal count-label budgets in this protocol. It remains information-unmatched: sparse process inputs contain all gold teacher-forced states.

\section{Verification and reproduction}\label{app:verify}\label{app:reproducibility}
\paragraph{Finite executors.}
\texttt{verification/lettertrace.py} preserves the Transformer and executor implementation. \texttt{verification/executor\_check.py} independently rebuilds the appendix tensors. The exhaustive driver enumerates all starts, queries and words for $n=2,\ldots,6$, alphabet size $2,3$, and $K=2,3,4$. It checks tensor agreement, every supervised next-token margin and complete greedy rollouts without teacher forcing in float64. All $31{,}635$ prompts and $363{,}654$ supervised predictions passed, with minimum observed margin exceeding $19.9995$. Results and minimum margins are in \texttt{evidence/executor\_verification.json}. This finite computation does not prove arbitrary-size correctness; Appendix~\ref{app:architectures} does.

\paragraph{Gradient identities.}
\texttt{verification/theorem4\_checks.py} tests the general-$K$ loss, derivatives, background symmetry, Frobenius gradient factors and elementary constants against dense autograd. It retains checks for the source misaligned branch, although this manuscript uses only aligned seeds. Numerical grids and finite differences detect errors but do not establish universal inequalities. The repository's Lean formalization contains unfinished proof obligations; the paper is not Lean-certified.

From \texttt{paper\_focused/}, run
\begin{verbatim}
make figures
make verify
make
\end{verbatim}
The build uses local TeX, styles, bibliography and figures. Figure regeneration reads parent-repository logs and writes compact evidence and SHA-256 provenance. Verification uses the copied implementation and requires no GPU. Override \texttt{PYTHON} or \texttt{TECTONIC} in the Makefile if needed. The parent \texttt{pyproject.toml} and \texttt{uv.lock} preserve software constraints.

Original pretrained experiments can be rerun from the repository root with model access and suitable compute:
\begin{verbatim}
python -m src lora --config configs/experiments/lora_transfer.yaml \
  --model Qwen/Qwen2.5-0.5B --output /tmp/qwen-main.csv \
  --run-dir /tmp/qwen-main-artifacts
python -m src lora --config configs/experiments/lora_transfer.yaml \
  --model Qwen/Qwen2.5-0.5B --conditions outcome --steps 6400 \
  --output /tmp/qwen-6400.csv --run-dir /tmp/qwen-6400
python -m src shuffled train --K 5 --condition process --seed 2003 \
  --out /tmp/bbh-K5-process-s2003.json
\end{verbatim}
Repeat BBH for $K\in\{3,5,7\}$, conditions \texttt{process,outcome}, seeds $2001,2002,2003$. Controlled counting uses \texttt{handcoded/regime\_both.py} with seed/stop/update environment variables documented there; label-control reproduction uses its original launch script. The reported model training was not rerun in this preparation.

\texttt{proof\_audit.md} records dependencies, metric correction, assumptions and unresolved qualifications; \texttt{claims\_matrix.md} links claims to proofs and evidence; \texttt{CHANGELOG.md} records source changes. These artifacts accompany the manuscript.

